\documentclass[11pt,a4paper]{article}
\usepackage[T1]{fontenc}
\usepackage[utf8]{inputenc}
\usepackage{lmodern,microtype}
\usepackage[margin=25mm,headheight=15pt]{geometry}
\usepackage{amsmath,amssymb,amsthm,mathtools,stmaryrd}
\usepackage{booktabs,tabularx,longtable,array}
\usepackage{xcolor,listings,enumitem,fancyhdr,titlesec,needspace,etoolbox}
\usepackage{xurl}
\usepackage[colorlinks=true,linkcolor=blue!38!black,citecolor=blue!38!black,urlcolor=blue!38!black]{hyperref}
\usepackage{bookmark}
\definecolor{ink}{HTML}{173348}
\definecolor{accent}{HTML}{21787A}
\definecolor{panel}{HTML}{F2F5F6}
\definecolor{muted}{HTML}{566773}
\titleformat{\section}{\Large\sffamily\bfseries\color{ink}}{\thesection}{.65em}{}
\titleformat{\subsection}{\large\sffamily\bfseries\color{ink}}{\thesubsection}{.65em}{}
\titleformat{\subsubsection}{\normalsize\bfseries}{\thesubsubsection}{.65em}{}
\pretocmd{\section}{\Needspace{10\baselineskip}}{}{}
\pretocmd{\subsection}{\Needspace{5\baselineskip}}{}{}
\setlength{\parindent}{0pt}
\setlength{\parskip}{.4em}
\setlength{\emergencystretch}{2em}
\setcounter{tocdepth}{1}
\setlist{nosep,leftmargin=1.6em}
\renewcommand{\arraystretch}{1.18}
\newcolumntype{Y}{>{\raggedright\arraybackslash}X}
\newcolumntype{P}[1]{>{\raggedright\arraybackslash}p{#1}}
\lstset{basicstyle=\small\ttfamily,columns=fullflexible,keepspaces=true,
 showstringspaces=false,breaklines=true,frame=single,framerule=.3pt,
 rulecolor=\color{accent!40},backgroundcolor=\color{panel},
 aboveskip=.7em,belowskip=.6em,tabsize=2}
\pagestyle{fancy}\fancyhf{}
\fancyhead[L]{\small\sffamily BASALT: properties, behavior, and mathematical reasoning}
\fancyhead[R]{\small\sffamily Round 3}
\fancyfoot[C]{\small\thepage}
\newcommand{\Basalt}{\textsc{Basalt}}
\newcommand{\code}[1]{\nolinkurl{#1}}
\newcommand{\R}{\mathbb R}\newcommand{\Q}{\mathbb Q}
\newcommand{\Z}{\mathbb Z}\newcommand{\N}{\mathbb N}
\newcommand{\eval}[1]{\llbracket #1\rrbracket}
\newcommand{\Spec}{\mathsf{Spec}}
\newcommand{\Pre}{\mathsf{Pre}}\newcommand{\Post}{\mathsf{Post}}
\newcommand{\Rfn}{\mathsf{Ref}}
\newcommand{\Obs}{\mathsf{Obs}}
\newcommand{\Req}{\mathsf{Request}}
\newcommand{\Fact}{\mathsf{Fact}}
\newcommand{\Perm}{\mathsf{Perm}}
\newcommand{\length}{\mathsf{length}}
\newcommand{\dom}{\mathsf{dom}}
\newcommand{\decision}[1]{\par\smallskip\textbf{Design decision.} #1\par\smallskip}
\theoremstyle{plain}
\newtheorem{proposition}{Proposition}[section]
\newtheorem{theorem}[proposition]{Theorem}
\newtheorem{lemma}[proposition]{Lemma}
\theoremstyle{definition}
\newtheorem{definition}[proposition]{Definition}
\theoremstyle{remark}
\newtheorem{remark}[proposition]{Remark}
\hypersetup{pdftitle={BASALT: A Property- and Behavior-Aware Mathematical Language over Lean},
 pdfauthor={OpenAI, prepared for Vladimir Reshetnikov},
 pdfsubject={Third-round proof-language proposal; refinement elaboration, relational contracts, certified computation},
 bookmarksnumbered=true}
\begin{document}
\begin{titlepage}
{\sffamily\color{accent}\large PROOF-LANGUAGE DESIGN\quad /\quad THIRD ITERATION\par}
\vspace{20mm}
{\sffamily\bfseries\color{ink}\fontsize{42}{46}\selectfont BASALT\par}
\vspace{7mm}
{\sffamily\color{ink}\fontsize{24}{29}\selectfont A Property- and Behavior-Aware\\Mathematical Language over Lean\par}
\vspace{9mm}
{\large Refinement-directed elaboration, relational contracts,\\scoped reasoning, and certified symbolic computation\par}
\vspace{12mm}
\rule{\textwidth}{.7pt}
\vspace{3mm}
\textbf{Central proposal.} Let a mathematician describe not only what kind of
object a term is, but which properties it has, which relations its maps preserve,
and what an operation promises to do. Infer the routine evidence needed to use
those descriptions, while emitting ordinary Lean terms and retaining the exact
statement, scope, and computational guarantee.
\vfill
Prepared for Vladimir Reshetnikov\par
6 September 2026\par
\vspace{4mm}
{\small\color{muted}\textbf{Evidence boundary.} Both round-two synthesis sources
were read closely, and selected Lean and Haskell modules were inspected directly.
This article specifies a language; it does not claim an implemented compiler or
measured usability improvement. The accompanying Python demonstrations were run.
The illustrative Lean core was not compiled in this environment.}
\end{titlepage}

\begin{abstract}
A mathematician rarely thinks of a function merely as a term of type
$A\to B$. It may be an increasing function on an interval, a measure-preserving
map, an ideal-preserving algebra homomorphism, or a procedure that returns all
solutions rather than one. Lean can already express these distinctions, but
ordinary authoring often separates the object from the proofs that make its
intended use legitimate. The resulting effort appears as repeated side-condition
assembly, coercion alignment, local-to-global bookkeeping, and transport between
representations.

\Basalt{} proposes a conservative extension of the \emph{elaborated type
language}: object-indexed refinements, dependent behavioral contracts, and
multi-object relational contracts become first-class guides to evidence
synthesis. The kernel's logic and definitional equality remain unchanged.
A local refinement enriches knowledge about the same term; an exported
refinement becomes a subtype or a record with proof fields. A relation such as
exactness belongs to an entire diagram, not to an adjective on each vertex.
Locality and precision are explicit indices of evidence, not informal levels
of confidence. An untrusted planner selects theorem applications, synthesis
providers, and certified computations; acceptance requires the exact requested
Lean proposition.

The proposal develops two substantial algebraic cases from the FLT repository,
a measure-theoretic case from ProveIt, and a locality control from its autonomous
ODE development. It also examines Leant's finite-spine length-contract handoff,
which distinguishes candidate provenance from behavioral proof. A practical
multivariate polynomial checker is specified down to coefficient denotation,
normalization, and original-target reconstruction. The article supplies typing
rules, composition theorems, failure examples, a conservative implementation
path, a fair evaluation design, and a source register. The proposed advance over
round two is not another name for a certified answer: it is a type-directed
account of how mathematical knowledge is accumulated and composed before an
answer is requested.
\end{abstract}

\tableofcontents
\clearpage

\section{The thesis: make mathematical knowledge part of the interface}
\label{sec:thesis}
The right target is not the shortest possible proof script. It is the shortest
\emph{faithful argument}: one that preserves the intended objects, assumptions,
quantifiers, and reasons while delegating details that can be reconstructed
reliably. A one-line invocation of a theorem identical to the goal is short, but
it does not explain a mathematical argument. Conversely, a long reusable lemma
may be the right place to pay for a difficult transport theorem once.

I propose \Basalt{}, a mathematical authoring layer whose fundamental judgment
is not simply ``this expression has this type,'' but:
\[
 \text{this exact object has this type and these proved properties,
 in this scope.}
\]
For maps and algorithms, the properties include \emph{behavior}: inputs meeting
specified conditions yield outputs satisfying specified relations, or related
inputs yield related outputs. For diagrams, properties can involve several
objects and arrows simultaneously. All these assertions elaborate to ordinary
Lean propositions and proofs.

The author should be able to write, schematically,
\par\noindent\begin{minipage}{\linewidth}
\begin{lstlisting}
let mu be an atomless probability measure on Real
let F := cdf mu
assume mu is invariant under (x |-> a - x)

claim F(a - x) = 1 - F(x)
  by pushforward, complement, and null-boundary replacement
\end{lstlisting}
\end{minipage}\par

The three named steps determine a calculational plan. The system retrieves the
probability and atomlessness evidence for this particular $\mu$, generates the
relevant measurable-set obligations, and checks the endpoint conversion. It
must not replace the plan with an unrelated theorem while continuing to claim
that those steps were verified. Nor may it silently add atomlessness when the
author supplied only a probability measure.

The design is intentionally hybrid. It combines a mathematical document,
a richer elaboration discipline, theorem-directed proof synthesis, and a
certificate-oriented CAS. None of these components gains authority from its
name. A mathematical property must have a proof, a computation must have its
advertised guarantee, and a suggestion must survive insertion at its original
target. This inherits the strongest boundaries of both second-round syntheses
\cite{synthesis,workbenches}.

\subsection{What is genuinely being extended}
Lean already has dependent functions, subtypes, structures with proof fields,
and propositions expressing arbitrarily rich mathematics. The proposed
extension is therefore \emph{not additional logical expressiveness}. It is a
new authoring and elaboration discipline that makes those existing facilities
work together with fewer explicit transports and less repeated hypothesis
assembly. In particular, it makes a distinction between discovering a proof
of a property, recording that property on a stable subject, and later using it
as a typing prerequisite.

An ordinary type can indeed be a behavioral specification when it is sufficiently
dependent or refined. The slogan ``a type is not a behavioral specification''
should be read as a warning about a \emph{coarse} type such as
$\mathsf{List}\,A\to\mathsf{List}\,A$, not as a limitation of dependent type
theory. A sorting contract, for example, can be a subtype of functions equipped
with a universal proof of sortedness and permutation preservation.

\subsection{Five commitments that distinguish this proposal}
The specific proposal consists of five interacting choices. First, locally
learning a property does not construct a new mathematical object. Second,
function contracts and diagram contracts are handled alongside unary
refinements. Third, evidence synthesis is driven by the requirements of the
current operation rather than by an attempt to infer every true property.
Fourth, locality, approximation, and representation change require typed
transport theorems. Fifth, abstract behavioral analysis is useful for pruning
and composition, but it becomes a Lean theorem only through proved primitive
laws and an interpretation theorem.

These are design choices to test, not claims of historical invention. Their
combination is developed against source examples in
Sections~\ref{sec:cdf}--\ref{sec:leant}, rather than justified by aesthetic
preference for the sample syntax.

\section{What the second round establishes, and what remains open}
\label{sec:roundtwo}
The two requested reports are \emph{Mathematical Objects, Certified
Computation} and \emph{Nine Workbenches}. Both were inspected as TeX sources,
including their corrections, evaluation proposals, and questions for a further
round. Their current versions explicitly discuss each other's revisions; they
should not be treated as untouched independent votes \cite{synthesis,workbenches}.

The first organizes agreement into eight broad commitments. The second gives a
finer editorial matrix and a consolidated mutation catalogue. The counts have
different denominators and do not establish either novelty or usability. Both
reports converge on an important acceptance boundary: a producer returns data
and evidence for the caller's \emph{fixed specification}. A conditional identity,
a finite jet, and a complete solution predicate are different mathematical
answers. A presentation remains connected to its original semantic object.

\begin{center}\small
\begin{tabularx}{\textwidth}{@{}P{33mm}Y Y@{}}
\toprule
Round-two principle & Retained in \Basalt & Third-round extension\\
\midrule
Fixed object and checked relation & No silent change of anchor or guarantee & Local refinements accumulate without replacing the anchor\\
Dependent guards & Conditions remain in their witness scope & Behavioral contracts derive operation-specific prerequisites\\
Annotated theorem interfaces & Ordinary Lean theorems are logical contracts & Refinement and relation indices select theorem roles\\
Exact-target providers & Acceptance is not a suggestive success label & Abstract behavioral laws must reconnect to exact Lean providers\\
Representation discipline & Preservation and reflection are separate & Commuting diagrams and observation consumers carry relational contracts\\
Retained evidence & Publication does not require rediscovery & Property dependencies support targeted invalidation and semantic merging\\
Fair Lean controls & Libraries and services are shared & Measure the incremental value of refinement-directed elaboration\\
\bottomrule
\end{tabularx}
\end{center}

\subsection{Corrections that must survive the next design}
Several repairs in the reports are especially instructive. A candidate solution
set needs both soundness and coverage; one witness plus coverage is insufficient.
Equality at a point does not become neighborhood equality merely because the
ambient domain is open. An enclosure can prove equality when its bounds
collapse, so a blanket prohibition would be too restrictive. A nonzero element
need not be a unit, and even a cancellable element need not be invertible.
Uniform sequence identities cannot be shortened by moving a quantifier inside
an equivalence \cite{synthesis,workbenches}.

These are not only checker problems. They are mistakes in choosing the
\emph{type of the requested evidence}. A stronger type discipline is valuable
precisely when it prevents these invalid promotions before search begins,
and explains which premise is missing when a proposed bridge does not apply.

\subsection{What not to repeat}
The earlier proposals already give elaborate versions of a dependent pair
$(y,p)$ with $p:\Spec(x,y)$. Repeating that interface under a tenth name would
add little. They also discuss the same binomial, autonomous-derivative, and
formal-series examples repeatedly. This article retains the autonomous example
as a control, but develops measure-theoretic and algebraic examples with
different proof obligations. Once inspected, these are \emph{development cases},
not held-out evaluation tasks.

Nor does this article reinterpret the reports' executed experiments as its own.
Their Lean probes and polynomial timings are historical evidence. No aggregate
ProveIt or FLT build was attempted here. The distinction between a proved
Boolean checker outcome and the unproved denotational soundness of that checker
is addressed explicitly in Section~\ref{sec:cas}.

\section{Fresh source diagnosis: four different kinds of verbosity}
\label{sec:diagnosis}
The source register in Appendix~\ref{app:sources} identifies the inspected
revisions and files. This is a focused design study, not a repository-wide
quality audit. The FLT repository's README reports extensive independent
checking; that report is not a checking result reproduced here. Its warning
that a declaration's statement, rather than its generated name or English
summary, determines what was proved is directly relevant to this project
\cite{fltreadme}.

\subsection{Repeated property assembly in measure theory}
ProveIt's \code{ContinuousCDF.lean} develops continuity of an atomless
probability CDF, reflection symmetry, and identification of a measure from an
everywhere derivative of its CDF. Its reflection calculation moves through
pushforward, preimage, complement, and an almost-everywhere equality of open
and closed half-lines. The type $\mathsf{Measure}\,\R$ alone does not tell the
elaborator which of these steps is available; the probability and
null-singleton hypotheses do \cite{cdfsource}.

The opportunity is not to hide those hypotheses. It is to attach their evidence
to $\mu$, derive the precise required consequences on demand, and retain the
reason for the endpoint step. In particular, the null-singleton property is
indexed by the measure: atomlessness of a different measure is irrelevant.

\subsection{Scope and locality in analytic induction}
The autonomous-derivative module proves an all-orders statement by induction.
The induction hypothesis holds on an open set, so at each point it supplies
eventual equality in the neighborhood filter. Only then may the proof replace
a derivative. The derivative assumptions on auxiliary functions are needed
only on the range of the original solution, not everywhere
\cite{odesource}.

This is not just positivity-style automation. It is a change in the kind and
scope of evidence: a universally quantified equality on an open set is
converted into local equality at a particular point. A property language that
records only ``the expressions are equal'' would erase the crucial distinction.

\subsection{Property-preserving maps in completions}
The FLT module \code{Def_AdicCompletionRingFunctoriality.lean} starts with an
algebra homomorphism $f:R\to S$ and ideals $I,J$, together with
$f(I)\subseteq J$. It builds maps on all quotients by powers of the ideals,
proves compatibility with transition maps, and lifts to completions. Many
lemmas repeatedly carry ideal-preservation evidence and relate the same map
through several bundled representations \cite{completionsource}.

Here the useful enriched type is not merely ``homomorphism.'' It is a morphism
between \emph{rings with selected ideals}, preserving the selected ideals.
Composition can then construct the preservation proof as part of its result.
The mathematical map remains $g\circ f$; the additional type information
explains when quotient and completion operations are legal.

\subsection{Relations on diagrams in exactness arguments}
The proof of \code{S_Rep_indBotSC_shortExact.lean} constructs a linear model of
an induced representation, proves that the modeling equivalences commute with
maps, tensors a short exact sequence with a flat module, and transports the
result back. Its final theorem is already short; the infrastructure establishing
the model and the commutative squares contains the substantive work
\cite{exactsource}.

Exactness is a property of a pair of composable arrows, and transport requires
a compatible ladder of equivalences. Attaching independent labels such as
``injective'' and ``surjective'' to isolated arrows does not capture all of
this. A language that only adds unary refinements will miss one of the main
sources of algebraic proof structure.

\decision{Treat unary properties, behavioral contracts, and diagram relations
as different surface forms of proved predicates, with different indexing and
elaboration policies. Do not force all mathematical knowledge into a hierarchy
of adjective-bearing object classes.}

\section{The surface language and its three views}
\label{sec:surface}
\Basalt{} has an authoring view, an expanded mathematical view, and an ordinary
Lean export. They are views of the same accepted evidence, not separate proofs
that are expected to agree by convention. All syntax in this article is
proposed unless explicitly identified as Lean.

\subsection{A small vocabulary}
The initial frontend needs a modest vocabulary: introduce objects; assert
properties; define a map by a contract; make a claim with a reason; calculate;
change representation; choose a witness satisfying a specification; and inspect
an unresolved obligation. Existing Lean terms remain usable inside these forms.

\par\noindent\begin{minipage}{\linewidth}
\begin{lstlisting}
context k : commutative ring
context (R, I), (S, J), (T, K) : k-algebras with ideals

let f : (R, I) => (S, J) preserving ideals
let g : (S, J) => (T, K) preserving ideals
let h := g after f

claim completion(h) = completion(g) after completion(f)
  by equality at every quotient level
\end{lstlisting}
\end{minipage}\par

The notation for pairs selects the actual algebra structures and ideals. It
does not ask typeclass search to invent an ideal, a scalar action, or a ring
structure. The final reason creates one schematic quotient-level equality,
not a finite test of several quotient levels.

\subsection{What appears in the expanded view}
For the preceding example, the expanded view displays $f,g$, the statements
$f(I)\subseteq J$ and $g(J)\subseteq K$, the derived statement
$(g\circ f)(I)\subseteq K$, the completion-map theorem used, and its
all-levels extensionality premise. Routine applications of membership lemmas
may be folded away. The choice of the scalar ring and the direction of the
maps may not.

A concise display is not necessarily a single line. If a solver returns a
conditional result, its condition belongs next to the result. If it returns
a finite observation, its precision belongs next to that observation. An
unresolved condition is not a cosmetic warning whose acknowledgment makes
it true.

\subsection{When ordinary Lean should remain the surface}
Short existing proofs should remain short. The final induction-of-exactness
corollary in the inspected FLT module already consists largely of applying
its prepared bridge. Translating it into a longer narrative would be a
regression. The first product should support lifting an existing Lean
statement into a property-aware display and selectively introducing new forms,
not demand conversion of an entire development.

\section{Extending the type system without extending the kernel}
\label{sec:types}
\subsection{What Lean already supplies}
Lean's subtypes represent values satisfying predicates, with a value field
and a proposition-valued proof field. They can be used like the underlying
value through a coercion, and their proof fields are computationally irrelevant
\cite{leansubtypes}. Bundled maps and ordinary structures already express many
mathematical interfaces. \Basalt{} should reuse these mechanisms rather than
replace them with a parallel mathematical universe.

The difficult part is systematic inference and use of those properties.
An author knows that the current measure is atomless; the next lemma expects
an almost-everywhere equality of two sets under that measure. A synthesis
procedure needs a route between these facts, must instantiate it at the right
subject, and should explain the route if it affects the argument.

\subsection{Three possible extension levels}
\textbf{Library only.} New bundled structures, theorem wrappers, and tactics
can expose the desired properties without changing syntax. This is the
necessary baseline and may be the best eventual product.

\textbf{Elaboration-level refinements.} A frontend accepts refinements,
behavioral arrows, and relational declarations, then emits ordinary Lean
terms with explicit proof arguments. This is the recommended experiment.
It changes the programmer-facing type system while preserving the
foundational one.

\textbf{Kernel-level refinements or conversion.} A primitive refinement
constructor could be given a conservative interpretation, but it would add a
new trusted implementation and migration surface for a feature already
representable by existing types. A more ambitious rule making theorem-level
property implication part of conversion would couple type checking to
arbitrary proof search. Neither change is justified as an initial step.

\decision{Implement an extended elaboration calculus first. Keep semantic
subtyping proof-producing and explicit in the elaborated output. Do not ask the
kernel to decide arbitrary predicate implication or to run a CAS as part of
definitional equality.}

\subsection{Local refinements versus exported values}
There are two distinct operations that should look related without being
confused.

\emph{Local enrichment} retains $x:A$ and a proof $p:P(x)$ in the existing
context. A property index records where that proof can be found. A later proof
$q:Q(x)$ adds knowledge about the same $x$. No rebinding to a new subtype value
is required.

\emph{Packaging} constructs a first-class value of
\[
 \Rfn(A,P):=\{x:A\mid P(x)\}.
\]
It is useful at module boundaries, in containers of certified objects, or when
the property is part of an API. It elaborates to a subtype or an ordinary
record. The two modes are linked by introduction and projection, but the
compiler must not introduce unnecessary wrappers after every local deduction.

For a row of properties $\mathcal P=(P_1,\ldots,P_r)$ on a fixed base type,
\[
 \Rfn(A,\mathcal P)=\{x:A\mid P_1(x)\land\cdots\land P_r(x)\}.
\]
This is an intersection of refinements of the \emph{same} object, not a product
of unrelated objects. A property such as preservation of a selected ideal
may mention parameters already in the context.

\subsection{Properties are not structure selection}
``$R$ is a field'' can mean either that a chosen ring structure has additional
laws or that one is introducing new data for a field structure. These are not
interchangeable. Multiple topologies, norms, scalar actions, or bases on the
same carrier can change a statement's meaning. A refinement may add a proof
about a selected structure; it must not silently select another structure.

The frontend therefore separates \code{with structure} from \code{satisfying
property}. Property closure can use a theorem about the current topology.
Changing the topology is a data-level operation with its own visible transport.
Likewise, a proof that an element is nonzero is not permission to replace a
commutative ring by a field.

\subsection{Why not turn every property into a typeclass?}
A small collection of familiar proposition-valued classes is useful for
interoperability. Nevertheless, turning every derived proposition into a global
instance would blur lexical scope, make search behavior difficult to predict,
and invite cycles between equivalent descriptions. A relational fact indexed
by a particular diagram is especially ill-suited to global instance search.

The proposed fact index is local, proof-bearing, and demand-driven. It may
install a local instance when a known Lean API requires one, but it retains the
proof's exact source and dependencies. Typeclass resolution for mathematical
structures and theorem search for property consequences remain distinct
operations with separately bounded policies.

\section{A small refinement-and-evidence calculus}
\label{sec:calculus}
This section specifies enough of the translation to make ``conservative'' a
checkable design claim. It is a metatheoretic design, not a completed
formalization of an elaborator.

\subsection{Contexts and judgments}
Let $\mathcal E$ be a Lean environment and $\Gamma$ an ordered dependent
context. Every accepted entry in a fact index has the form
\[
  (\text{subject tuple }\vec t,\ \varphi,\ p),\qquad
  \mathcal E;\Gamma\vdash p:\varphi(\vec t).
\]
The fact index stores handles and provenance; it does not enlarge the logic.
Its entries must already be derivable in the Lean context, or be explicit
hypotheses in that context. A displayed subject name is not its identity.

Write $e\rightsquigarrow t:A$ for elaboration of a source expression to a
Lean term, and $e\Leftarrow T\rightsquigarrow t$ for checking against an
expected enriched type. During elaboration, obligations may be unresolved;
then the result is a \emph{pending derivation}, not an accepted term.
Publication requires all proof holes and data holes to be resolved under the
chosen axiom policy.

The initial enriched type grammar is deliberately small:
\[
 T ::= A\ \mid\ \{x:A\mid\varphi(x)\}\ \mid\ (x:T)\to T'
       \ \mid\ \text{an ordinary dependent record}.
\]
Behavioral and diagram contracts elaborate into instances of this grammar.
The mathematical predicate language is Lean's proposition language, not a
second restricted logic pretending to express every theorem.

\subsection{Refinement introduction and elimination}
The introduction rule is
\[
\frac{\Gamma\vdash e\rightsquigarrow t:A\qquad
      \Gamma\vdash p:\varphi(t)}
 {\Gamma\vdash e\Leftarrow\{x:A\mid\varphi(x)\}
       \rightsquigarrow\langle t,p\rangle}.
\]
Its premises show the separation of term elaboration from evidence synthesis.
An obligation solver can propose $p$, but it cannot change the already fixed
meaning of $t$ to make $\varphi(t)$ easier.

Elimination projects the value and makes its proof available:
\[
 z:\{x:A\mid\varphi(x)\}
 \quad\leadsto\quad z.\mathsf{val}:A,
 \qquad z.\mathsf{property}:\varphi(z.\mathsf{val}).
\]
A local enrichment is just a local proof declaration about $t$. A row
projection uses conjunction elimination; adding a property uses conjunction
introduction. None of these rules needs a new axiom.

\subsection{Subtyping means a checked adapter}
Given a proof
\[
 h:\forall x:A,\ P(x)\to Q(x),
\]
there is a value-preserving adapter
\[
 \mathsf{weaken}_h:\Rfn(A,P)\to\Rfn(A,Q),\qquad
 \langle x,p\rangle\mapsto\langle x,h(x,p)\rangle.
\]
The elaborator may insert it automatically at a checking boundary. The
insertion is not a new definitional equality between the two subtype types.
Even when $P\leftrightarrow Q$, a propositional equivalence is translated into
maps or an explicit transport, not silently installed as kernel conversion.

Proof irrelevance helps with different proofs of the \emph{same} proposition.
It does not identify different values, different predicates, different chosen
structures, or incompatible relation instances. The data identity must remain
visible in every adapter.

\subsection{Dependent application and unresolved guards}
Suppose $f:(x:\Rfn(A,P))\to B(x.\mathsf{val})$ and $a:A$. Applying $f$ to
$a$ generates the obligation $P(a)$ and elaborates, once that obligation is
proved, to $f\langle a,p\rangle$. If $P(a)$ remains open, the application is
pending. A frontend may offer a conditional theorem instead, but only as a
different result whose condition is shown.

When outputs or later premises depend on previous witnesses, the order of the
telescope is retained. The forms
\[
 \forall h:G,\ \{y:B\mid R(h,y)\}
 \quad\text{and}\quad
 \{y:B\mid\forall h:G,\ R(h,y)\}
\]
are different interfaces. No binder may be moved simply because its name is
hidden in the compact display. If the type of $y$ itself depends on $h$, even
the second expression needs a different formulation.

\subsection{Relative soundness of the core rules}
\begin{proposition}[Conservative translation of the displayed core]
Assume the base expressions and registered theorem constants are well typed in
$\mathcal E$, and every proof premise supplied to the displayed rules is checked
by Lean in its indicated context. A finite accepted derivation using base
elaboration, refinement introduction and projection, conjunction rules,
checked weakening, dependent abstraction and application, and explicit equality
transport translates to a well-typed term of its interpreted type in Lean.
\end{proposition}
\begin{proof}
Induct on the derivation. Base cases are the assumed Lean judgments.
Refinement introduction is the subtype constructor with its two checked
arguments; projection is a structure projection. Row rules are constructors
or projections of conjunctions. Weakening applies the supplied implication to
the value and its existing proof. Abstraction and application are the
corresponding dependent function rules. Equality transport applies Lean's
ordinary equality eliminator. The induction hypothesis supplies each
subderivation at the exact context required by the next rule.
\end{proof}

This proposition is deliberately modest. It explains why the suggested rules
need no new logic. It does not verify a parser, a mutable elaboration engine,
a serialization format, the property index, or the default display. Those
components need implementation tests and, where warranted, formal verification.
It also gives no completeness theorem for automatic evidence search.

\section{Behavioral types: functions with mathematical obligations}
\label{sec:behavior}
A useful function interface says more than which objects can be passed to it.
It can state which preconditions suffice, what is guaranteed about the output,
and which relations the function preserves. These are properties of the
\emph{selected function}, not facts guessed from its name.

\subsection{A total function with a conditional contract}
For $f:A\to B$, a deterministic input--output contract is
\[
 \mathsf{Contract}(f,P,Q)
 :=\forall x:A,\ P(x)\to Q(x,f(x)).
\]
The predicate $Q$ can refer to both input and output. A packaged function is
\[
 \{f:A\to B\mid\mathsf{Contract}(f,P,Q)\}.
\]
Here $f$ remains a total Lean function; the contract promises the specified
behavior only under $P$. This is different from a function whose domain is
itself $\Rfn(A,P)$, and different again from a partial computation returning
an optional result. The surface may support all three, but its default must
preserve the meaning of an existing Lean declaration.

\par\noindent\begin{minipage}{\linewidth}
\begin{lstlisting}
construct transform : List A -> List A
  ensures (xs, ys) => ys is a permutation of xs

construct sort : List A -> List A
  ensures (xs, ys) => sorted(ys) and permutation(ys, xs)
\end{lstlisting}
\end{minipage}\par

The first contract permits reversal, identity, and many other functions. The
second excludes identity on unsorted inputs. Neither contract, by itself,
requires a particular sorting algorithm, stable ordering of elements with
equal keys, or a running-time bound. Such demands need additional predicates
or an explicitly modeled execution semantics.

\subsection{Composition derives a real proof, not a label}
Let $f:A\to B$ and $g:B\to C$. Suppose
\begin{align*}
 h_f&:\forall x,\ P(x)\to Q(x,f(x)),\\
 h_g&:\forall x\,y,\ Q(x,y)\to R(x,y,g(y)).
\end{align*}
Then their composition has the contract
\[
 \forall x,\ P(x)\to R(x,f(x),g(f(x))).
\]
Its proof is the term
\[
 \lambda x\,h_P.\ h_g\,x\,(f(x))\,(h_f\,x\,h_P).
\]
The intermediate output is retained. This matters when $g$'s precondition
refers to the result selected by $f$, rather than to an arbitrary object with
similar printed syntax.

A common specialization has an intermediate invariant $Q(y)$ independent of
$x$. A more general sequential contract introduces a witness together with
its relation to $x$. Neither requires an untyped bag of guard strings.
Backward planning can derive a sufficient precondition for the whole sequence,
but \Basalt{} does not promise to compute weakest preconditions for arbitrary
mathematical predicates.

\subsection{Variance and consequence}
Suppose $f$ has a contract with precondition $P$ and postcondition $Q$. To use
it at a new interface $(P',Q')$, it suffices to prove
\[
 \forall x,\ P'(x)\to P(x),\qquad
 \forall x\,y,\ P'(x)\to Q(x,y)\to Q'(x,y).
\]
The new interface may demand a stronger precondition, but it cannot silently
claim the old implementation works for more inputs. It may advertise a weaker
postcondition, but not a stronger one without a theorem. For dependent
functions, the output adapter is instantiated after the input and its proof
are fixed.

This directionality prevents a common cache mistake. A theorem established
under $x\ne0$ is not a stronger result than an unconditional theorem merely
because it has an additional hypothesis. A higher-precision observation and a
result with stronger assumptions vary in different directions.

\subsection{Relational behavior and abstraction adequacy}
Some properties compare two executions:
\[
 \mathsf{Preserves}(f,\mathcal R,\mathcal S)
 :=\forall x\,x',\ \mathcal R(x,x')\to\mathcal S(f(x),f(x')).
\]
Examples include monotonicity, preservation of congruence modulo an ideal,
nonexpansiveness, and compatibility with an observational equivalence. These
contracts permit certified replacement under selected consumers without
claiming the underlying objects are equal.

A useful general bridge is the following.
\begin{proposition}[Adequacy of an observation]
Let $o:A\to O$ be an observation, $f:A\to B$ an operation, and
$\widehat f:O\to B$ an abstract implementation. If
$\forall x,\ f(x)=\widehat f(o(x))$, then $o(x)=o(y)$ implies $f(x)=f(y)$.
\end{proposition}
\begin{proof}
Substitute the two adequacy equations and apply congruence of $\widehat f$ to
$o(x)=o(y)$.
\end{proof}
For example, equality of list lengths is adequate for a function known to
factor through length. It is not adequate for the identity function on lists,
a sorting result, or a sum of elements. Rather than declaring an observation
``safe,'' a package proves which consumers respect it. This supplies a
concrete interpretation of observational types without inventing a universal
hierarchy of information.

\subsection{Universal behavior is not finite testing or polymorphism}
Testing a candidate on finitely many lists proves, at most, the tested cases.
A transfer function for lengths becomes universally valid only after the
semantics of the candidate's permitted constructors and primitives have been
connected to that transfer function by proof.

Nor should a Haskell-inspired synthesizer import a global parametricity
principle into Lean. Lean's official axiom discussion gives noncomputable
polymorphic functions that inspect a type proposition using classical
reasoning, invalidating a naive naturality theorem for every polymorphic
list function \cite{leanaxioms}. A restricted syntax can support a proved
parametricity or length theorem; an unrestricted Lean type does not provide
that theorem automatically.

This is an especially important distinction for Leant. Structural synthesis
can remain excellent at constructing candidate programs. Its abstraction's
behavioral laws need an explicit semantic interpretation before they become
properties of the corresponding Lean functions.

\section{Relational and scoped types: what belongs to the whole situation}
\label{sec:relations}
\subsection{Beyond a row of adjectives}
A short exact sequence, a commutative square, a measure pushforward identity,
and a function with a specified inverse are predicates on several subjects.
They should be introduced as named relation instances:
\par\noindent\begin{minipage}{\linewidth}
\begin{lstlisting}
let S := (A --f--> B --g--> C)
know S is short exact

let e := (eA, eB, eC)
know e is an isomorphism of diagrams from S to T
\end{lstlisting}
\end{minipage}\par

The first declaration does not follow from $f$ being injective and $g$ being
surjective alone. The middle exactness condition remains necessary. The second
requires commuting squares; three arbitrary objectwise isomorphisms are not
enough. The compiled relation is an ordinary Lean proposition or a record
whose proof fields state the required equations.

The fact index therefore uses subject tuples or a canonical diagram handle,
not only one object identifier. It is a hypergraph of proved relationships in
an implementation sense, but no graph algorithm is granted mathematical
authority. Every edge contains, or reconstructs, a Lean proof.

\subsection{Locality is an index, not a confidence score}
The following are distinct predicates:
\begin{align*}
 &f=g,\\
 &\forall x\in U,\ f(x)=g(x),\\
 &f=^{\mathcal F}g\quad\text{(eventual equality in a specified filter)},\\
 &f=g\quad\mu\text{-almost everywhere},\\
 &[t^k]f=[t^k]g\quad\text{for }k<N.
\end{align*}
The last expression concerns formal series rather than arbitrary functions.
These predicates do not form one linear order. In particular, almost-everywhere
equality is usually not enough to preserve a value or derivative at a
specified point. A finite jet supplies exact finite algebraic information,
not a numerical error estimate.

Some transfers are legitimate and useful. If $V\subseteq U$, equality on $U$
restricts to $V$. Equality on an open neighborhood of $x$ yields eventual
equality in $\mathcal N(x)$. Almost-everywhere equality under $\mu$ transports
to $\nu$ when a suitable absolute-continuity theorem applies. Each transfer
has its own direction and premises. A mere inclusion of printed domain labels
is not evidence for the transfer.

\subsection{Intersection, images, and dependent locality}
Composing two on-domain equalities produces equality on the intersection of
the domains. A request on a larger domain is not thereby solved. Likewise,
using a property of $g$ on $V$ while composing $g\circ f$ on $U$ requires
$f(U)\subseteq V$. This image-containment fact is relational and depends on
the selected $f$, $U$, and $V$.

For differentiation, a neighborhood statement about a composition may need
continuity or a derivative theorem to establish the relevant image behavior.
The autonomous example avoids a gratuitous global strengthening by asking its
auxiliary derivative hypotheses only at values $W(x)$ for $x\in U$.
The initial frontend should preserve such theorem signatures rather than
replace them with an easier but stronger ``every function is smooth globally''
profile.

\subsection{A common shape without a universal relation lattice}
A package may define an evidence family
\[
 \Obs(\vec x;\iota):\mathsf{Prop},
\]
where $\vec x$ are its subjects and $\iota$ are indices such as a set, filter,
measure, precision, branch, or selected diagram. It registers concrete
weakening and composition theorems. This gives common infrastructure for
search, storage, and rendering while leaving mathematical meaning in the
package's actual predicates.

A root-isolating interval and a complete algebraic solution predicate can be
incomparable. Two valid interval enclosures need not be equal as intervals.
Two exact presentations of the same object can be related through their
meaning proofs without forcing every representation route to normalize to a
single expression. Coherence is required only for the observation the caller
actually uses.

\section{Evidence synthesis: demand-driven, bounded, and subject-preserving}
\label{sec:synthesis}
\subsection{A concrete elaboration procedure}
The intended procedure has two phases. First, elaborate the statement and
meaning-bearing data: variable types, scalar structures, domains, relation
indices, and named methods. Second, generate and solve the proof obligations
required by the chosen expression or method. An author may explicitly delegate
a data choice by a fixed specification; otherwise proof search must not assign
meaning-bearing holes merely to make the goal easy.

For each required proposition, the initial implementation tries direct local
evidence, projections from accepted refined values, and explicitly registered
small implication rules. It then tries bounded theorem application, selected
ordinary Lean tactics, and optional providers such as Leant or a certificate
producer. The order and budgets are project policy, not mathematics.

\par\noindent\begin{minipage}{\linewidth}
\begin{lstlisting}
elaborate application at fixed statement and revision:
  infer the selected operation and its dependent roles
  instantiate its ordinary Lean contract
  for each premise, in telescope order:
    reuse exact local evidence, if available
    otherwise request evidence under a bounded profile
    validate the returned term at this exact premise
  build the application from the accepted arguments
  check the result at the original requested type
  commit the source edit only if the revision still matches
\end{lstlisting}
\end{minipage}\par

This describes a procedure, not a completeness claim. Failure reports the
unsolved premise or exhausted route. It does not mean that the theorem is false.

\subsection{Index rules by the conclusion's subjects and mathematical role}
A rule for continuity of a CDF should be discoverable from a requirement about
$\mathsf{cdf}(\mu)$ and refer back to the same $\mu$. A rule for composition of
ideal-preserving maps should identify its source ideal, intermediate ideal,
and target ideal. A rule for exactness transport should identify both arrows
and the square witnesses.

These roles are annotations on stable project-owned theorem wrappers. They
are not inferred indefinitely from fragile binder positions or from English
names. The underlying theorem remains callable in ordinary Lean. A wrapper
signature change is checked as an interface change, not silently accepted
because a similarly named declaration still exists.

\subsection{Why not saturate the context with every consequence?}
Unrestricted forward closure is both unbounded and undesirable. From one
property, there may be infinitely many consequences and many equivalent
formulations. Aggressive expansion also obscures which assumptions a fact
actually uses.

The preferred implementation performs demand-driven backward search with
memoization, supplemented by a very small collection of cheap forward rules.
A memo entry is indexed by the elaborated proposition, its ordered context,
relevant environment identity, and policy. A proof term, not a Boolean
``known'' flag, is the authority. Exact identity is the cheap reuse path;
reuse across changed contexts needs checked instantiation and any required
transport.

\subsection{Guards, cycles, and ordinary induction}
A conditional theorem $G\to Q$ can be recorded even when $G$ is not known.
It cannot be used to establish $Q$ while simultaneously using that same
unproved $Q$ to establish $G$. A guard solver detects such unresolved cycles
and reports them rather than manufacturing evidence.

This does not outlaw mathematical induction or recursive definitions.
An induction hypothesis is a legitimate scoped premise supplied by an
induction theorem or a well-founded recursion principle. The distinction is
between a checked rule authorizing recursive reasoning and a circular chain
of unresolved elaboration obligations with no such rule.

\subsection{Reuse and maintenance}
A derived property records the terms and hypotheses it actually depends on.
Changing a measure invalidates properties about that measure; renaming its
printed binder need not. Changing a selected topology may invalidate analytic
properties even when the carrier has the same name. Changing a proof of an
unchanged proposition need not change the mathematical subject.

For cross-revision reuse, hashes identify candidates for reuse, not semantic
proofs. An adapter must check that the current target follows from retained
evidence under the current hypotheses and policy. The dependency representation
should begin as an index into Lean's existing context and declaration data,
not as a second, independently authoritative environment.

\subsection{Diagnostics should name the mathematical obstruction}
A useful failure is not ``cannot synthesize instance 37.'' It is, for example:
\begin{quote}
The null-boundary step requires that this measure assigns zero mass to
$\{x\}$. The context contains the probability and reflection laws, but no
proof of that boundary condition. The calculation is valid up to the
open-half-line expression; the claimed closed-half-line equality remains open.
\end{quote}
This message can be generated from typed roles and the unresolved proposition.
It need not be a free-form explanation invented after the fact. Expanded views
should expose the exact Lean target behind every such description.

\section{Worked case I: probability, atomlessness, and CDF symmetry}
\label{sec:cdf}
Let $\mu$ be a probability measure on $\R$, with no atoms, and let
$r_a(t)=a-t$. Suppose $(r_a)_*\mu=\mu$. Write
$F(x)=\mu(({-}\infty,x])$. The ProveIt theorem
\code{cdf_reflection_sub} establishes
\[
 F(a-x)=1-F(x)
 \qquad\text{for every }x\in\R
\]
under these hypotheses \cite{cdfsource}.

\subsection{A complete mathematical calculation}
The reflection map is measurable, and
$r_a^{-1}(({-}\infty,a-x])=[x,\infty)$. Therefore
\begin{align*}
 F(a-x)
 &=\mu(({-}\infty,a-x])\\
 &=((r_a)_*\mu)(({-}\infty,a-x])\\
 &=\mu([x,\infty))\\
 &=1-\mu(({-}\infty,x))\\
 &=1-\mu(({-}\infty,x])\\
 &=1-F(x).
\end{align*}
The fourth equality uses total mass one and measurability. The fifth uses
$\mu(\{x\})=0$. These are different proof obligations, even though a human
usually reads them without pausing.

\Basalt{} would register a probability refinement on $\mu$, an atomless
refinement on $\mu$, and a reflection-invariance relation involving $r_a$
and $\mu$. The calculation requests the appropriate consequences, rather
than expanding every property at declaration time.

\par\noindent\begin{minipage}{\linewidth}
\begin{lstlisting}
let mu : Measure Real satisfying probability and atomless
let F := cdf mu
assume pushforward(reflection(a), mu) = mu

calculate F(a - x)
  = mass(mu, [x, infinity))       by pushforward reflection(a)
  = 1 - mass(mu, (-infinity, x)) by probability complement
  = 1 - F(x)                    by null-boundary replacement
\end{lstlisting}
\end{minipage}\par

The first displayed step abbreviates the definitional bridge from CDF to
half-line mass. Its retained evidence must include that bridge and the preimage
identity. This is exactly where statement-oriented code and calculation-oriented
presentation need a common source-to-evidence map.

\subsection{A negative neighbor: reflection without atomlessness}
Take $\mu=\delta_0$ and $a=0$. This is a reflection-invariant probability
measure, but $F(0)=1$. The proposed symmetry at zero would assert $1=0$.
The missing null-boundary proof is essential, not bureaucratic overhead.

A well-designed interface can still return the valid statement
\[
 F(a-x)=1-\mu(({-}\infty,x)),
\]
or equivalently the correct closed-half-line formula with the atom correction
\[
 F(a-x)=1-F(x)+\mu(\{x\}).
\]
For the latter, decompose $({-}\infty,x]$ into its open part and singleton.
The original stronger request is not discharged unless the correction term
is shown to vanish. This is a useful positive alternative, rather than a
system that merely refuses every difficult step.

\subsection{A second theorem: derivative behavior determines the measure}
The same source proves that if a function $f$ satisfies
\[
 \forall x\in\R,\quad \mathsf{HasDerivAt}(F,f(x),x),
\]
then $\mu$ is Lebesgue measure with density $f$, expressed in Lean by
\code{withDensity} using $\mathsf{ENNReal.ofReal}(f(x))$
\cite{cdfsource}. Its proof uses monotonicity of the CDF to derive
$f\ge0$ and interval integrability of the derivative, applies the fundamental
theorem of calculus on each interval, and identifies measures on half-open
intervals.

An enriched declaration can express the whole derivative relation:
\par\noindent\begin{minipage}{\linewidth}
\begin{lstlisting}
let f : Real -> Real
know derivative(F) is f at every real point

claim mu = lebesgue.withDensity(f)
  by agreement on half-open intervals using the fundamental theorem
\end{lstlisting}
\end{minipage}\par

The phrase ``at every real point'' belongs in the public contract. It cannot
be dropped in favor of ``almost everywhere.'' The step CDF of $\delta_0$ has
derivative zero away from zero, hence almost everywhere, but $\delta_0$ is
not the zero-density measure. Continuity or differentiability properties
must therefore carry their exact quantifier and locality indices.

\subsection{Where the conciseness comes from}
This example does not need an exotic new prover. It needs theorem packages
for CDF observations, local evidence about the selected measure, and a
calculation interface that keeps endpoint and quantifier distinctions visible.
The gain would come from not reassembling the same contextual prerequisites
at every step. Whether that gain exceeds what good Lean wrappers alone provide
is an empirical question reserved for the matched evaluation.

\section{Worked case II: ideal-preserving maps and completion functoriality}
\label{sec:completion}
Fix a commutative scalar ring $k$. Consider commutative $k$-algebras $R,S,T$
with selected ideals $I,J,K$. The inspected FLT source constructs the maps
below explicitly, including their quotient-level compatibility and their
composition laws \cite{completionsource}.

\subsection{A useful dependent morphism type}
Define the interface
\[
 \mathsf{IdealHom}_k((R,I),(S,J))
 =\{f:R\to_{k\text{-alg}}S\mid f(I)\subseteq J\}.
\]
The notation $f(I)\subseteq J$ can be implemented by the existing ideal-map
inclusion predicate. The selected ideals are part of the type's parameters.
The interface does not bundle an arbitrary new ideal with each use of $f$.

\begin{proposition}[Composition of ideal-preserving maps]
If $f(I)\subseteq J$ and $g(J)\subseteq K$, then
$(g\circ f)(I)\subseteq K$.
\end{proposition}
\begin{proof}
For $r\in I$, the first premise gives $f(r)\in J$ and the second gives
$g(f(r))\in K$. The composite is the ordinary composite algebra homomorphism.
\end{proof}
Consequently composition returns a value of the enriched morphism type,
with its proof field obtained by composition of the two preservation proofs.
This is a small behavioral theorem, not a new algebraic operation.

An initial wrapper can expose this proof in ordinary Lean. The new syntax is
justified only if it makes repeated use of the wrapper easier to author,
understand, or maintain.

\subsection{From ideal preservation to maps on quotient towers}
For every $n$, the homomorphism sends $I^n$ into $J^n$: products of $n$
elements of $I$ map to products of $n$ elements of $J$, and the generated
ideal is respected. Thus it induces
\[
 f_n:R/I^n\longrightarrow S/J^n,\qquad [r]\longmapsto[f(r)].
\]
The preservation condition is exactly what makes this representative formula
well defined. For $m\le n$, let $q^R_{nm}$ and $q^S_{nm}$ be the quotient
transition maps. Then
\[
 q^S_{nm}\circ f_n=f_m\circ q^R_{nm},
\]
since both sides take the class of $r$ to the class of $f(r)$. The source
proves corresponding identities by passing to a representative and reducing
the definitions \cite{completionsource}.

These equations form a \emph{relation on the whole family} $(f_n)_n$.
They should be retained as a tower-morphism contract, not rediscovered from
unrelated local facts at each use. A map between the carrier types of two
levels, even if bijective, is not automatically a morphism of towers.

\subsection{Completion as a contract-preserving construction}
In the compatible-system description of the completion, an element is a family
$x_n\in R/I^n$ satisfying the transition equations. Applying $f_n$ to each
coordinate preserves those equations, so it defines
\[
 \widehat f:\widehat R_I\longrightarrow\widehat S_J.
\]
The source implements the construction with its existing completion lifting
interface and evaluation maps, rather than with a new definition of completion
\cite{completionsource}. A \Basalt{} package should wrap that interface and
its characterizing equations.

\begin{proposition}[Functoriality by observations]
For ideal-preserving algebra maps $f$ and $g$,
\[
 \widehat{g\circ f}=\widehat g\circ\widehat f.
\]
\end{proposition}
\begin{proof}
At each quotient level $n$, both sides send the coordinate represented by
$r$ to the coordinate represented by $g(f(r))$. The quotient-level composition
law gives equality of their evaluations for every $n$. Extensionality of
completion elements then gives equality of the maps.
\end{proof}
The phrase ``for every $n$'' is essential. Checking a finite prefix of the
tower supplies a finite observation, not the functor law.

\par\noindent\begin{minipage}{\linewidth}
\begin{lstlisting}
let f : IdealHom(k, (R,I), (S,J))
let g : IdealHom(k, (S,J), (T,K))

claim complete(g after f) = complete(g) after complete(f)
proof by all quotient observations
  fix n
  both maps send [r] to [g(f(r))]
end
\end{lstlisting}
\end{minipage}\par

The final sentence is an abbreviation for quotient induction and the
characterizing equations. Those are the method contract. An untrusted planner
can fill the term details, but cannot replace the universal observation with
sample computations.

\subsection{Levelwise equivalence needs compatibility}
The same source constructs a completion equivalence when the induced map is
bijective at every level. Compatibility of the inverse maps follows from
compatibility of the forward maps and injectivity at the lower level. To see
this, apply $f_m$ to the proposed equation
\[
 q^R_{nm}f_n^{-1}(y)=f_m^{-1}q^S_{nm}(y).
\]
Naturality makes both images equal to $q^S_{nm}(y)$, and injectivity of $f_m$
finishes the proof. The inverse coordinate family is therefore compatible.

When bijectivity is supplied only propositionally, constructing an inverse
may use noncomputable choice. The source explicitly uses noncomputable
definitions. A contract promising an equivalence should not thereby advertise
an executable inverse algorithm. Conversely, a caller requesting an executable
inverse must supply or synthesize an appropriate computational presentation.

\subsection{Negative neighbors and their typing failures}
For the identity ring map on $\Z$, choose $I=(2)$ and $J=(3)$. The condition
$I\subseteq J$ fails. The representative formula from $\Z/(2)$ to
$\Z/(3)$ would send the equal classes of $0$ and $2$ to unequal classes.
The enriched morphism type correctly prevents constructing this quotient map.

A forward preservation condition for an algebra equivalence does not imply the
reverse preservation condition for its inverse. Similarly, arbitrary
objectwise equivalences of quotient levels need not commute with transitions.
Neither omission is repaired by displaying the word ``equivalence.''

The desirable gain is the automatic composition of \emph{available} behavioral
proofs. It is not automatic discovery of an entire completion theorem from a
coarse function type.

\section{Worked case III: exactness is a relational type}
\label{sec:exactness}
The inspected FLT proof models an induced representation using a tensor
product with a free module and proves that this model is compatible with maps.
It then derives preservation of short exactness
\cite{exactsource}. This is a strong test because it combines chosen
representations, mathematical properties of modules, and relations on diagrams.

\subsection{The abstract ladder argument}
Let
\[
 0\longrightarrow M_1\xrightarrow{f}M_2\xrightarrow{g}M_3
 \longrightarrow0
\]
be a short exact sequence of $k$-modules and let $F$ be a flat $k$-module.
Suppose another diagram $S$ has objects $S_1,S_2,S_3$ and maps $f',g'$,
and there are linear equivalences
\[
 e_i:S_i\overset{\sim}{\longrightarrow}F\otimes_k M_i
 \quad(i=1,2,3)
\]
such that
\begin{align*}
 e_2\circ f'&=(1_F\otimes f)\circ e_1,\\
 e_3\circ g'&=(1_F\otimes g)\circ e_2.
\end{align*}
The sequence obtained by tensoring is short exact because $F$ is flat. The
equivalences and commuting squares transport its three components: left
injectivity, middle exactness, and right surjectivity.

For middle exactness, take $b\in S_2$ with $g'(b)=0$. Commutativity puts
$e_2(b)$ in the kernel of $1_F\otimes g$. Tensor exactness supplies $u$ with
$(1_F\otimes f)(u)=e_2(b)$. Surjectivity of $e_1$ gives $a$ with $e_1(a)=u$.
The first square and injectivity of $e_2$ yield $f'(a)=b$. The reverse inclusion
follows from the zero-composition law. Injectivity and surjectivity are
transported by analogous applications of the squares and the equivalences.

This is the mathematical content of the source's ladder helper. The actual
application uses the free module of finitely supported scalar functions on
the group, so the flatness requirement is discharged by the relevant module
structure rather than by strengthening the base ring to a field
\cite{exactsource}.

\subsection{A diagram-oriented interface}
\par\noindent\begin{minipage}{\linewidth}
\begin{lstlisting}
let X be a short exact sequence of k-representations
let U := underlying_module_sequence(X)
let F := finitely_supported_functions(G, k)
know F is flat                     by its free-module structure

let T := tensor(F, U)
know T is short exact              by flat tensor transport

let e := induced_representation_model(X)
know e is a diagram isomorphism from underlying(indBot(X)) to T
  by naturality on generators

conclude indBot(X) is short exact  by diagram transport
\end{lstlisting}
\end{minipage}\par

The word ``naturality'' generates the two required squares. A method for
checking them on generators needs a spanning or induction principle; it is
not permission to test a few sample vectors. The source explicitly develops
such a generator-based induction and its formulas for the modeling
isomorphism \cite{exactsource}.

The relevant property types are therefore relational:
\[
 \mathsf{ShortExact}(f,g),\qquad
 \mathsf{Commutes}(e_1,e_2,f,f'),\qquad
 \mathsf{DiagramIso}(S,T,e).
\]
Their dependencies make the larger proof readable. A flatness fact about an
unrelated module or over a different scalar ring does not satisfy the tensor
step.

\subsection{What a unary system would lose}
Even for a short complex, injectivity of the first arrow and surjectivity of
the second are insufficient. Over a nonzero field, consider
\[
 k\longrightarrow k^3\longrightarrow k,
 \qquad a\longmapsto(a,0,0),\quad (x,y,z)\longmapsto z.
\]
The maps compose to zero and have the two endpoint properties, but the
kernel at the middle is two-dimensional while the image is one-dimensional.
A missing middle-exactness relation is a genuine mathematical gap.

Flatness also matters. Multiplication by $2$ is injective on $\Z$, but after
tensoring with $\Z/2\Z$, the corresponding endomorphism is zero on a nonzero
module and is not injective. A generic ``tensor preserves properties'' label
would therefore be wrong. The exact preservation theorem and its hypotheses
must determine what the language infers.

\subsection{An honest compression account}
The final source theorem already uses its prepared ladder helper compactly.
The new language should not claim to save the helper's mathematical work by
hiding it behind one keyword. Its plausible contribution is to make the
helper's reusable interface explicit, to discover the correct instances and
squares, and to expose missing relational prerequisites before low-level
coercion mismatches dominate the error message.

\section{Control case: range-local analytic induction}
\label{sec:ode}
The autonomous-derivative module is deliberately retained as a control from
the earlier rounds. Let $U$ be open, $W'=\phi\circ W$ on $U$, and suppose a
family $G_n$ has derivatives $G'_n$ at all values $W(x)$ for $x\in U$, with
\[
 G_0(W(x))=W(x),\qquad
 G_{n+1}(W(x))=G'_n(W(x))\phi(W(x)).
\]
Then $W^{(n)}(x)=G_n(W(x))$ on $U$. This is the contract of the inspected
source theorem, with derivative existence represented explicitly
\cite{odesource}.

A proposed proof reads:
\par\noindent\begin{minipage}{\linewidth}
\begin{lstlisting}
claim iterated_derivative(n, W) = G(n) after W on U
proof by induction on n, retaining the domain U
  zero: use the initial identity
  successor:
    differentiate the previous identity locally on U
    use the chain rule at values attained by W
    use the recurrence for G
end
\end{lstlisting}
\end{minipage}\par

The induction hypothesis has an on-$U$ index. At $x\in U$, openness permits
its conversion to eventual equality in $\mathcal N(x)$. The local derivative
transfer theorem then applies, followed by the chain rule and the recurrence.
The range-local derivative contract is sufficient; global differentiability
of every $G_n$ would be an unnecessary strengthening.

The negative neighbor is $f(t)=t$ and $g(t)=0$ at $t=0$. Equality of their
values at zero cannot justify replacing their derivatives. A second negative
neighbor removes derivative existence and keeps only a displayed equation
using Lean's total derivative operator; such an equation does not provide the
\code{HasDerivAt} premise needed by the chain rule. The system must request
the stronger evidence rather than silently reinterpret the notation.

This example tests whether the refinement layer improves use of an already
well-factored theorem. It does not establish that new type syntax is necessary:
an annotated Lean wrapper may perform equally well.

\section{Leant integration: from structural candidates to proved behavior}
\label{sec:leant}
\subsection{What was inspected, and what the code actually promises}
The inspected Leant revision is
\code{3a40904be8a410d832d9ab6900b3d3e7b425eccb}, the same revision cited by the
round-two reports. The additional finding here is not a later version; it is
functionality beyond the narrow verification and tactic-suggestion discussion
highlighted in the syntheses.

The verification module's opaque \code{Verified} wrapper records acceptance
by a supplied callback. It explicitly disclaims behavioral evidence, a solver
certificate, and a kernel proof. The engine exposes exact typed origins and
retained semantic sidecars. The length-contract vocabulary describes passive
assertions, including explicitly supplied provider laws. The length handoff
then binds one accepted candidate occurrence to a candidate-specific
finite-spine problem using its retained origin and exact rendering
\cite{leantverification,leantengine,leantcontract,leanthandoff}.

This is useful existing infrastructure. It should not be dismissed as merely
``a string suggester,'' nor promoted to an original-target behavioral theorem.
The handoff's stated guarantee is checked structural correspondence at its
particular boundary. Its refusal cases include changed targets, missing typed
origins, renderer changes, and unavailable or ambiguous provider identities.
Those distinctions should survive the frontend integration.

\subsection{The missing bridge for theorem-level behavioral claims}
An asserted provider law such as ``this provider preserves length'' is not
a proof about the corresponding Lean function. A \Basalt{} adapter should
require a theorem witness indexed by that exact function. For a primitive
$u$, an example law is
\[
 h_u:\forall xs,\ \length(u(xs))=L_u(\length(xs)).
\]
The adapter then needs a theorem relating the permitted candidate syntax to
its denotation and a theorem that the abstract length interpretation matches
that denotation. Finally it must connect the denoted candidate to the actual
Lean expression inserted at the target.

The chain is
\[
 \begin{gathered}
 \text{exact candidate origin}\longrightarrow\text{typed candidate syntax}\\
 \longrightarrow\text{Lean denotation}\longrightarrow\text{behavioral theorem}.
 \end{gathered}
\]
The first arrow may be an integration invariant supported by structured
receipts; the latter semantic implications require actual proofs before the
whole result can be advertised as a Lean theorem. Naming every stage
``checked'' would conceal precisely the distinction the source preserves.

\subsection{A restricted length calculus with an explicit soundness argument}
Consider a finite first-order grammar
\[
 e ::= x_i\mid []\mid\mathsf{reverse}(e)\mid e_1\mathbin{+\!+}e_2,
 \qquad 1\le i\le r.
\]
Let $\eval e(\vec{xs})$ be its ordinary list semantics. Define an abstract
interpretation $L(e)\in\N^r$ by
\[
 L(x_i)=\mathbf e_i,\quad L([])=0,\quad
 L(\mathsf{reverse}(e))=L(e),\quad L(e_1\mathbin{+\!+}e_2)=L(e_1)+L(e_2).
\]
\begin{proposition}[Length interpretation]
For every expression in this grammar and every list input vector,
\[
 \length(\eval e(\vec{xs}))
 =\sum_{i=1}^{r}L(e)_i\,\length(xs_i).
\]
\end{proposition}
\begin{proof}
Induct on $e$. The input case selects the corresponding length, and the empty
list has length zero. Reversal preserves length. In the append case, the
length of the result is the sum of the two lengths, so substitute the two
induction hypotheses and distribute the finite sum.
\end{proof}
The mathematical proof identifies exactly which primitive laws a Lean
formalization needs. Extending the grammar with another primitive, recursion,
or case analysis requires extending the denotation and the proof; a
recognizable function name is not a substitute.

The companion implements this restricted interpreter and tests finite
instances. It is not an implementation of the full Leant candidate language,
and the proposition above has not been kernel-checked in this work.

\subsection{A subtle negative-evidence issue: realizable abstract inputs}
Even an exact length calculation does not automatically make every abstract
counterexample a concrete input. For a fixed element type $A$, the set
\[
 V_A=\{\length(xs):xs:\mathsf{List}\,A\}
\]
is $\{0\}$ when $A$ is empty, and all of $\N$ when $A$ is inhabited. An
arithmetic counterexample requiring an input of length one cannot refute a
claim about lists over an empty type.

Thus a counterexample provider needs both an abstract violating assignment
and a \emph{realization} by inputs of the original types. A concrete element
of $A$ can build lists of the required lengths; a mere unsupported assumption
that such an element exists cannot. Positive abstract proofs can remain sound
without every abstract input being realizable, but negative evidence requires
this additional bridge.

This is an important place for dependent behavioral types to help. The
counterexample result should contain a concrete input and a proof that it
violates the candidate's contract, rather than a detached vector of lengths.
A refutation of one candidate is also not a refutation of the existence of
some other candidate satisfying the synthesis request.

\subsection{Correlated outputs and relation strength}
The inspected source also contains a binary-product finite-spine handoff
\cite{leanthandoff,leantcontract}. For a pair of output lists, a behavioral
contract should preserve correlations involving the same input, rather than
store two unrelated ``length-like'' labels. For example,
\[
 \length(ys)+\length(zs)=\length(xs)
\]
is different from two independent upper bounds. The interpretation theorem
must match the exact product representation and source-ordered roles.

Length alone does not prove permutation, content preservation, sorting, or
identity. Both reversal and identity preserve length, and identity fails a
sorting request on $[2,1]$. These are not defects in a length domain; they are
invalid promotions between observations. The frontend should make the
accepted behavioral predicate explicit.

\subsection{One provider envelope, several mathematical authorities}
A shared provider envelope can carry the request identity, ordered context,
source revision, selected policy, proposed term or certificate, and diagnostic
status. It need not impose one universal certificate language. A proof term,
a polynomial certificate, a typed length derivation, and a counterexample
have different semantic checking theorems.

For tactic edits, the actual displayed edit must be replayed in its original
surrounding context after insertion. For data synthesis, the selected value
and its specification must be checked. For negative outcomes, distinguish
unsupported input, bounded failure, nonderivability in a stated calculus,
refuted candidate, and a checked negation of the original target. This is a
proposed adapter contract; this article does not claim to have re-audited every
current tactic-suggestion path in \code{Main.hs}.

\section{Certified CAS: a concrete multivariate checker, not a slogan}
\label{sec:cas}
A hybrid system should use computation to discharge mathematical obligations
whose guarantees fit the current refined type. The first verified core need
not be a complete CAS. A small polynomial certificate checker is enough to
exercise coefficient domains, reification, execution, and exact-target
reconstruction.

\subsection{Finite representation and denotation}
Fix a number of variables $d$. A monomial exponent is a vector
$\alpha\in\N^d$. Represent an integer polynomial by a finite coefficient map
$c:\N^d\to\Z$, stored as a sorted list of nonzero coefficients with unique
exponent vectors. The zero polynomial has an empty support. A parser verifies
that every exponent vector has length $d$, every exponent is a natural number,
and every coefficient is an integer.

For a commutative ring $R$ and a valuation $v:\{1,\ldots,d\}\to R$, define
\[
 \eval c_v=\sum_{\alpha\in\operatorname{supp}(c)}
      \iota_R(c_\alpha)\prod_{j=1}^{d}v_j^{\alpha_j},
\]
where $\iota_R:\Z\to R$ is the integer scalar map. No injectivity of
$\iota_R$ is assumed. Empty products equal one, so the definition also handles
$d=0$.

Addition merges equal exponent vectors and adds their coefficients.
Multiplication forms all pairs of terms, adds exponent vectors, multiplies
coefficients, and normalizes the resulting list. Normalization groups duplicate
exponent vectors and deletes zero coefficients. These are concrete finite
algorithms, not operations on an unspecified semantic polynomial type.

\subsection{Why the arithmetic operations preserve denotation}
\begin{lemma}[Normalization]
Grouping terms with the same exponent vector, replacing their coefficients by
the sum, and removing zero coefficients does not change denotation.
\end{lemma}
\begin{proof}
For one exponent vector $\alpha$, distributivity and additivity of $\iota_R$
give
\[
 \sum_i\iota_R(a_i)\prod_jv_j^{\alpha_j}
 =\iota_R\!\left(\sum_i a_i\right)\prod_jv_j^{\alpha_j}.
\]
Regrouping a finite sum in a commutative additive group preserves its value,
and zero terms contribute zero. Apply this independently to each group.
\end{proof}

\begin{lemma}[Addition and convolution]
For the finite operations just described,
\[
 \eval{c+c'}_v=\eval c_v+\eval{c'}_v,
 \qquad
 \eval{c*c'}_v=\eval c_v\,\eval{c'}_v.
\]
\end{lemma}
\begin{proof}
Addition follows from the normalization lemma. For multiplication, expand the
product of the two finite sums. Multiplicativity of $\iota_R$ and the power
law in a commutative ring identify a pair of monomial terms with the monomial
whose exponent vector is their sum. Normalize the resulting finite list using
the preceding lemma.
\end{proof}
These proofs require commutative multiplication. A noncommutative expression
must use a different representation and theorem, not a permissive coercion to
this checker.

\subsection{The certificate and its soundness theorem}
Given polynomials $p,f_1,\ldots,f_m$, a certificate is a vector of polynomials
$q_1,\ldots,q_m$. The checker verifies the exact coefficient identity
\[
 p=\sum_{i=1}^{m}q_i f_i.
\]
The arity $m$ must agree. A parser must not silently truncate mismatched input
lists with a zip operation. A typed implementation can use vectors indexed
by the same finite type after validating the external representation.

\begin{theorem}[Denotational soundness of the specified checker]
If the described checker accepts $(p,\vec f,\vec q)$, then for every
commutative ring $R$ and valuation $v$,
\[
 \eval p_v=\sum_{i=1}^{m}\eval{q_i}_v\,\eval{f_i}_v.
\]
In particular, if $\eval{f_i}_v=0$ for every $i$, then $\eval p_v=0$.
\end{theorem}
\begin{proof}
Acceptance means that the two canonical coefficient maps are identical.
Identical maps have identical denotations by the definition of the finite
sum. Repeatedly apply the addition and convolution lemmas to the right-hand
map. Under the vanishing hypotheses, each summand is zero.
\end{proof}
This is a universal mathematical soundness proof of the specified algorithm.
It is not a claim that the Python implementation is verified or that this
proof has been accepted by Lean. A Lean implementation must formalize the
finite representation and these lemmas, connect its Boolean equality test to
coefficient-map equality, and check the actual computation.

\subsection{The original-target bridge}
Suppose the caller's goal concerns an expression $e:R$ and hypotheses
$h_i:e_i=0$. Reification must produce $p,f_i$ and proofs
\[
 \rho:e=\eval p_v,\qquad \rho_i:e_i=\eval{f_i}_v.
\]
Given a checked acceptance proof and the soundness theorem, rewrite the
hypotheses through $\rho_i$, obtain $\eval p_v=0$, and transport through
$\rho$ to conclude $e=0$.

Opaque subexpressions may become variables in $v$, but their identities must
be exact elaborated terms in the original context. Two atoms with the same
printed name in sibling scopes are not interchangeable. If two syntactically
different atoms are identified by a theorem, the reification bridge must
include that theorem. Coefficient-domain casts require the appropriate
homomorphism laws, not a heuristic string conversion.

A suitable Lean contract has the following shape; the names are illustrative:
\par\noindent\begin{minipage}{\linewidth}
\begin{lstlisting}
check_sound :
  check input certificate = true ->
  forall (R : Type) [CommRing R] (v : Fin d -> R),
    denote target v = sum_i (denote multiplier[i] v *
                              denote generator[i] v)
\end{lstlisting}
\end{minipage}\par

The consumer supplies the reification equalities and actual vanishing
hypotheses. The checker does not need to understand the caller's whole theorem.

\subsection{Execution is another proof obligation}
A correctness theorem about a Lean function does not certify bytes returned by
an unchecked executable. The actual assertion
$\mathsf{check}(i,c)=\mathsf{true}$ must be obtained through the chosen trusted
execution route: kernel reduction, checked reconstruction, or an explicitly
approved computation mechanism. Lean's validation documentation separates
statement interpretation, axioms, and proof checking; publication must account
for each \cite{leanvalidation}.

Native evaluation is not authorized merely because it is faster. A project's
allowed axiom inventory and execution assumptions are a prior policy choice;
performance selects among the routes that policy permits. Toolchain upgrades
require a new audit of the transitive dependencies, rather than renaming an
old allowlist entry. The source reports' native-versus-kernel timings should
not be generalized to a different checker or machine \cite{synthesis,workbenches}.

\subsection{Limits that the type of the result should expose}
The certificate proves an ideal-membership consequence. Failure to find a
certificate is not nonmembership. A certificate for $p^n$ does not in a general
ring prove $p=0$: $2^2=0$ in $\Z/4\Z$, but $2\ne0$. Such a promotion needs an
appropriate reducedness or power-reflection theorem.

A coefficient extension may preserve an identity without reflecting a richer
existence claim. For example, $X$ belongs to the ideal generated by $2X$ in
$\Q[X]$, using multiplier $1/2$, but not in $\Z[X]$. A rational certificate
cannot be read as an integer-coefficient witness. Likewise, equality of
polynomial functions over a finite field need not be equality of polynomials:
$X^2-X$ is nonzero as a polynomial over $\mathbf F_2$ but vanishes at every
field element. The requested relation determines which checker is appropriate
\cite{synthesis,workbenches}.

\subsection{What the companion actually demonstrates}
The Python companion implements the sparse integer representation, its
normalization and convolution, exact-arity certificate checking, and the
restricted length interpreter of Section~\ref{sec:leant}. Its 24 test methods
passed in this environment. They include valid multivariate certificates,
corrupted coefficients, mismatched arities, malformed exponents, a finite-jet
residual, and concrete counterexamples for invalid property promotions.

These tests are finite implementation demonstrations. They do not formalize
reification of Lean expressions, validate Leant's actual candidate semantics,
measure Lean checker performance, or establish usability. The universal
arguments in this section are written mathematical proofs; the missing
kernel formalization remains an explicit first implementation deliverable.

\section{Observation-indexed contracts connect jets and quotient towers}
\label{sec:precision}
The formal-series proposals in round two already developed demand-driven
precision planning. The completion example suggests a wider application of
the same idea: make precision behavior a relational contract on operations,
not a special heuristic inside a CAS \cite{synthesis,workbenches}.

\subsection{A general sufficient-demand interface}
Suppose each object type $A$ has registered observation relations
$E_A(n,x,y)$, meaning that $x$ and $y$ agree at the observation level $n$.
For an operation $f:A\to B$, a sufficient-demand transformer is a function
$d_f:\N\to\N$ with a proof
\[
 \forall n\,x\,y,\quad
 E_A(d_f(n),x,y)\to E_B(n,f(x),f(y)).
\]
The interface says how much agreement at the input suffices for a requested
agreement at the output. It need not be optimal. Monotonicity or restriction
rules for observations are additional theorems when needed by a planner.

\begin{proposition}[Composition of precision contracts]
If $f:A\to B$ has sufficient-demand transformer $d_f$ and $g:B\to C$ has
$d_g$, then $g\circ f$ has transformer
\[
 d_{g\circ f}(n)=d_f(d_g(n)).
\]
\end{proposition}
\begin{proof}
Input agreement at $d_f(d_g(n))$ gives agreement of the two $f$-outputs at
$d_g(n)$ by the contract for $f$. The contract for $g$ then gives output
agreement at $n$.
\end{proof}
The order of composition is important: the consumer's demand propagates
backward through $g$ before $f$. For multi-input operations the transformer
returns a vector of sufficient demands, with any guards and shared parameters
retained in their actual scope.

\subsection{Three concrete instances}
For formal series, let $E(n,A,B)$ mean agreement of coefficients below degree
$n$. The formal derivative has sufficient demand $d_D(n)=n+1$, since the
coefficient of degree $k$ after differentiation uses the input coefficient of
degree $k+1$. This is a formal algebraic statement, not a theorem about the
analytic derivative of an arbitrary function.

For the ideal-preserving maps in Section~\ref{sec:completion}, let $E_R(n,x,y)$
mean $x-y\in I^n$ and similarly for $S,J$. The induced map has demand
$d_f(n)=n$, by preservation of powers of ideals. Its extension to completion
observations satisfies the corresponding quotient-level contract.

For truncated polynomial multiplication, agreement below degree $n$ in each
input suffices for agreement below degree $n$ in the product. More precise
rules can use known valuations, but they require their own proved hypotheses.
The planner must not infer a valuation merely from a coefficient list that
was never connected to the original series.

\subsection{Residual certificates remain observations of the original root}
A particularly useful inherited bridge starts with an existing formal root
$\Delta$ of a polynomial equation $F(Y)=0$ over $R[[t]]$. If $J$ has the same
constant term, the reduced derivative at that term is a unit, and
$F(J)$ vanishes below degree $N\ge1$, then $J$ and $\Delta$ agree below degree
$N$. Indeed,
\[
 F(J)-F(\Delta)=(J-\Delta)U,
\]
where the constant term of $U$ is the specified unit. Thus $U$ is invertible,
and membership of the residual in $(t^N)$ implies membership of
$J-\Delta$ in $(t^N)$ \cite{synthesis,workbenches}.

A computed $J$ is an observation of the original $\Delta$, not a replacement
for a universally quantified root. The result's type records its precision
and branch data. A demand for the first $N$ coefficients of $D\Delta$ can
then request $N+1$ coefficients of $\Delta$, using the derivative contract.
The finite producer still needs a proof that its truncation algorithm realizes
the advertised observation.

This does not collapse formal series, completion elements, analytic germs,
and numerical approximations into one type. It provides a shared protocol for
relation-indexed information flow. Every domain supplies its own interpretation,
restriction theorems, and operation contracts.

\section{Definitions, choices, and the boundary between mathematics and execution}
\label{sec:definitions}
\subsection{Definition packages should expose reasons for existence}
A mathematical definition package can contain a semantic object, an existence
or uniqueness theorem, characterization lemmas, observations, and certified
computational presentations. The package need not expose every implementation
detail in the default mathematical view, but it must not introduce a hidden
axiom when a construction is missing.

There are three different author actions:
\par\noindent\begin{minipage}{\linewidth}
\begin{lstlisting}
fix Delta satisfying its defining equation
construct Delta satisfying its defining equation
choose Delta from a proved existence theorem
\end{lstlisting}
\end{minipage}\par

The first introduces a quantified object or an explicit hypothesis. The
second generates an existence-and-construction obligation. The third uses a
particular existence theorem, possibly with noncomputable choice. The compiler
must retain these distinctions. A canonical constructed solution cannot simply
replace the arbitrary object introduced by the first form without a uniqueness
or observation theorem connecting them.

\subsection{Choice by specification need not require constant prompting}
Automation may choose a witness when the author explicitly asks for any witness
satisfying a fixed specification. It need not ask permission for every candidate
formula. What it may not do is change that specification, narrow the domain,
choose a different branch, or drop a required uniqueness or completeness clause.

For example, constructing a nonnegative real square root needs both the square
equation and nonnegativity. The negative root has the same square but fails
the specified branch. Constructing all roots of a polynomial requires a
complete solution predicate or another representation with soundness and
coverage; one certified root is insufficient.

\subsection{Preserve total operations when lifting Lean}
Existing Lean statements use their existing total operations. A frontend must
not quietly replace total real division by partial-expression semantics. For
example, in the total interpretation,
\[
 \frac{x^2-1}{x-1}=x+1
\]
fails at $x=1$, where the left side is zero and the right side is two.
A guarded equality for $x\ne1$, a correct piecewise identity, and an equality
of partial expressions on a domain are three distinct possible requests.

An opt-in partial-expression package may use a domain predicate and a value
on that domain. Its simplification contracts must preserve or explicitly
change the domain as well as the values. The source view should visibly select
this package; it should not reinterpret an imported Lean expression merely
because the partial reading is more familiar to a mathematician.

\subsection{Behavioral constraints beyond output equations}
A proof about a pure total function does not specify the performance of every
implementation with that extensional behavior. A complexity claim needs a
chosen cost semantics, for example an instrumented evaluation returning
$(y,c)$ together with a proof of $c\le B(x)$. A memory or state invariant
likewise needs a state-transition model and pre/postconditions for that model.
The refinement machinery can express these contracts, but cannot infer their
semantics from an ordinary function type.

External search may time out or fail even when the mathematical function being
specified is total. A timeout is an operational outcome, not a partial proof
or a mathematical counterexample. Noncomputable mathematical choice is also
different from a failed execution: it can define an object with a checked
specification without providing an executable algorithm for its values.

\decision{Use behavioral types for mathematical input--output and relational
specifications first. Add cost or state effects only together with an explicit
execution semantics. Do not let a generic ``verified'' annotation blur these
different claims.}

\section{Three correctness boundaries, not one green checkmark}
\label{sec:boundaries}
\subsection{Theorem truth}
The first boundary is the Lean proposition actually proved, in its actual
environment and with its actual axiom dependencies. Refinements, guards,
behavioral laws, and transport proofs all contribute to this boundary. The
small-core translation argument in Section~\ref{sec:calculus} explains why
these features can remain conservative.

Inconsistent explicit hypotheses still permit arbitrary conclusions, as in
ordinary logic. The frontend does not solve mathematical consistency for every
context. It must instead show which hypotheses the theorem assumes and never
silently promote a desired conclusion into an assumption.

\subsection{Statement and source fidelity}
The second boundary is whether the elaborated proposition matches the source
that the author meant to submit. A perfectly checked proof of a changed
statement is not a solution to the original task. Meaning-bearing metavariables,
structure selection, overloaded operations, quantifier dependencies, and
branch conventions therefore need explicit treatment before proof search.

A statement seal is a readable expansion of the actual interpretation, not
a certificate of human intention. The default view must show restrictions
that change the claim. The fully expanded Lean statement remains available.
The FLT repository's own distinction between declaration names and authoritative
statements is a useful reminder that presentation fidelity cannot be inferred
from kernel validity alone \cite{fltreadme}.

\subsection{Method and workflow fidelity}
The third boundary is whether the accepted artifact performs the method or edit
it claims to perform. A proof may establish the right theorem without using
the named argument. A tactic probe may succeed while the displayed text fails
when inserted into the document. A result may be valid but stale after the
user changes a witness or scope.

A named method therefore generates a structured derivation record containing
its subclaims and the theorem instances linking them. Publication checks those
subclaims, including unused displayed assertions. Because Lean treats proofs
of a proposition irrelevantly, the proposition alone does not encode a unique
explanation. Method fidelity is a validation policy on the elaborated derivation
and source map, not a new axiom that proofs reveal their historical origin.

\subsection{An honest status representation}
A result should distinguish at least mathematical guarantee, proof-validation
state, unresolved guards, source binding, and policy compatibility. For example,
a theorem $G\to Q$ can be fully checked while $G$ remains unproved at the
caller. A finite observation can be fully checked while not answering an exact
equality request. A proof can be valid at an old source revision while no
longer eligible to modify the current file.

These distinctions need not appear as a cumbersome vector in ordinary reading.
They should determine the wording and available actions. ``Checked conditional
result; boundary condition unresolved'' is better than a green badge followed
by a footnote that contradicts it.

\section{Architecture and implementation plan}
\label{sec:implementation}
\subsection{Lean owns the authoritative objects}
The implementation should retain Lean expressions, local contexts, declaration
identities, and source locations as authoritative inputs. An additional index
maps subject tuples and proposition heads to existing proof handles. It also
records method roles, selected interpretation packages, and retained computation
certificates. It should not duplicate Lean's entire context management.

The main components are a frontend parser, an elaboration extension, a
proof-bearing property index, a registry of ordinary theorem contracts,
provider adapters, and an evidence-aware renderer. The finite checker lives
in the Lean library once verified. An external CAS, large-language-model
planner, or Haskell synthesizer is untrusted with respect to theorem truth.

\begin{center}\small
\begin{tabularx}{\textwidth}{@{}P{37mm}Y Y@{}}
\toprule
Component & Existing foundation & New obligation\\
\midrule
Local objects and facts & Lean contexts and proof terms & Subject/role index must not retarget evidence\\
Mathematical structures & Existing bundled structures and instances & Make selected structures visible; reject ambiguous data choices\\
Property inference & Theorem application and tactics & Bounded rule policies and exact-subject reconstruction\\
Diagram transport & Existing exactness and naturality theorems & Checked role mappings and coherent diagram packaging\\
CAS acceptance & Proved checker plus execution proof & Exact reification and original-target reconstruction\\
Leant adapter & Typed origins and callback receipts & Proved semantic laws and behavioral interpretation\\
Document rendering & Source spans and elaborated information & Display the actual scope and strength of every assertion\\
Publication and migration & Ordinary Lean checking and dependency audit & Search-free retained evidence and new receipts after upgrades\\
\bottomrule
\end{tabularx}
\end{center}

\subsection{A minimal complete slice}
The first implementation should not attempt all the syntax in this article.
It should support local property enrichment, one refined morphism wrapper,
one relation on a diagram, and one certified computation request. Each must
also be usable directly from ordinary Lean.

A useful first milestone is the following combination. Formalize the sparse
polynomial checker and its denotational soundness, provide exact reification
for a restricted ring-expression grammar, and use it in guarded cancellation.
In parallel, package ideal-preserving algebra maps and prove their composition
wrapper using the existing completion APIs. These two tasks distinguish
computational certificates from theorem-level property propagation while
sharing the same typed request boundary.

The first frontend commands can simply be \code{know}, a refined function
application, and a named calculation step. A narrow default display can show
the required premises and their sources. A complete prose grammar, an
unrestricted planner, and a global representation search engine are not
prerequisites.

\subsection{Subsequent exit gates, not speculative schedules}
The next gate is a CDF calculation whose probability and boundary evidence is
derived from the correct measure, with the Dirac negative case diagnosed at
the missing boundary condition. A third gate is the exactness ladder, with
commuting squares and flatness kept explicit. A fourth connects the Leant
length domain to proved primitive laws and the exact accepted Lean candidate.

Each gate requires a positive case, a nearby negative case, retained evidence,
and successful checking with discovery services disabled. A source edit or
library upgrade must be replayed before the gate is considered passed again.
The repository snapshots use different toolchains; an implementation must pin
one environment and check any ports explicitly rather than assume their
compiled artifacts are interchangeable.

\subsection{Safe transactional elaboration}
Candidate attempts run against an immutable request snapshot. Shared
metavariables are committed only after the candidate and all dependent
obligations have been validated. A failed attempt restores the prior state.
A late result is rejected when the originating statement or source revision
has changed, even if its proof was valid in the old context.

Operational commit tokens are not mathematical propositions to be freely
assumed. They are application-level permissions tied to a revision. There is
no need to impose linear logic on ordinary mathematical facts merely to
prevent stale editor writes. The security and workflow layer should remain
separate from the reusable mathematics library.

\subsection{Publication artifacts and collaboration}
Publication retains source, its elaborated statements, theorem terms or
checkable certificates, the relevant environment pin, and a map from displayed
steps to evidence. Rediscovery scripts may be useful for editing but should
not be the only retained authority. External services disappearing must not
invalidate already retained mathematics.

When two branches of a document are merged, merge data choices and proof
evidence separately. Two new proofs of an unchanged property usually coexist.
Two changed witnesses or two conflicting representation routes require a new
semantic check. A union of fact caches is not a sound merge procedure when
those facts belong to different source contexts.

\section{Relation to prior type-system and proof-language work}
\label{sec:prior}
The proposal should not present refinements themselves as new. Lovas and
Pfenning develop refinement sorts for logical frameworks and interpret them
through proof irrelevance, including a bidirectional discipline. Ghalayini
and Krishnaswami develop explicit refinement types with proof terms rather
than relying only on a fixed SMT fragment, and report a Lean formalization of
their system \cite{lovaspfenning,explicitrefinements}. These works support the
choice to make evidence explicit while keeping user-facing annotations compact.
They are precedents, not proofs of the correctness of \Basalt.

Lean's existing subtype mechanism supplies an immediate representation for
value refinements, while ordinary dependent records and functions supply the
behavioral and relational encodings. The novelty claim here is therefore
limited: a proposed integration discipline specialized to mathematical
workflows, developed against the cited source examples and the present Leant
handoff. No claim is made that the underlying refinement calculus, Hoare-style
consequence rule, certificate checking, or functoriality arguments are newly
discovered mathematics.

Relative to the nine round-two proposals, the emphasis changes from the type
of a \emph{returned computation} to the type-directed accumulation and use of
knowledge \emph{throughout} an argument. The most important additions are
relational diagram contracts, contextual enrichment without object replacement,
proved behavioral-abstraction bridges, countermodel realizability, and the
application of observation-demand contracts to both symbolic and algebraic
examples. These are refinements of an already strong foundation, not a
replacement for its trust boundaries \cite{synthesis,workbenches}.

\section{Evaluation: make the new layer earn its place}
\label{sec:evaluation}
\subsection{Separate library, service, and language effects}
The primary study should compare four matched conditions: existing Lean;
Lean with the new theorem packages; Lean with the same packages and providers;
and \Basalt{} with exactly those packages and providers. An evidence-aware
rendered document view can be tested separately. Otherwise a language could
claim credit for a theorem that was simply absent from the baseline.

All wrapper authoring, checker proofs, role annotations, test maintenance,
and training effort count toward the cost. Record unsuccessful attempts as
well as completed proofs. Match solver budgets and library access. A theorem
used while designing a package is development material, even if its filename
had not appeared in an earlier report.

\subsection{A concrete workload portfolio}
The development portfolio comprises the CDF reflection and derivative-density
tasks, ideal-preserving completion maps, the short-exactness ladder, the
range-local derivative control, guarded polynomial reasoning, and restricted
list behavior. It spans measure-indexed properties, unary map constraints,
multi-arrow relations, filter-local evidence, and semantic abstraction.

Independent maintainers should select genuinely held-out tasks only after the
interfaces are frozen, excluding aliases, direct restatements, and near-duplicate
downstream consequences. Some held-out tasks should deliberately remain short
in ordinary Lean. The frontend should have a low-cost escape hatch rather
than require unnecessary narrative scaffolding.

\subsection{What to measure}
The main outcome is time and effort to a correctly stated, checked, maintainable
theorem, not visible line count alone. Secondary outcomes include author
interventions, wrapper reuse, unresolved-obligation quality, proof-checking
cost, certificate size, and maintenance effort after controlled changes.
Distinguish discovery time, elaboration time, and kernel checking time.

A semantic-reading test asks readers to reconstruct the exact subjects,
assumptions, quantifier scope, domain, branch, precision, and completeness of a
compact statement before inspecting its expansion. A method-reading test asks
which inference justified a displayed transition and which facts it required.
A concise interface that systematically induces wrong answers to those
questions has failed, even if it saves keystrokes.

For Leant, classify obligations into direct lookup, structural composition,
behavioral checking, and original-target reconstruction. Measure it on actual
obligations generated by the packages rather than assuming its general
synthesis ability automatically helps every CAS request.

\subsection{Paired semantic acceptance tests}
The following is an implementation backlog. Only the explicitly identified
finite instances in the companion were executed; this table is not a claim
that a language implementation passed the suite.

{\small
\begin{longtable}{@{}P{29mm}P{52mm}P{64mm}@{}}
\toprule
Boundary & Positive case & Nearby negative or required alternative\\
\midrule\endfirsthead
\toprule Boundary & Positive case & Nearby negative or required alternative\\\midrule\endhead
\bottomrule\endfoot
Boundary mass & Atomless reflection gives CDF symmetry & Dirac reflection needs the atom correction; unconditional symmetry fails\\
Derivative scope & Everywhere CDF derivative supplies the source theorem's premise & Almost-everywhere derivative alone does not identify a density\\
Subject identity & Atomlessness evidence for the selected measure & Evidence for another measure is not reusable by matching its printed name\\
Quotient map & Ideal-preserving homomorphism induces every level map & Identity from the $(2)$ quotient to the $(3)$ quotient is not well defined\\
Map composition & Composed ideal-preservation proof & Missing intermediate containment leaves a guard open\\
Diagram transport & Objectwise equivalences and commuting squares & Equivalences alone do not establish a diagram isomorphism\\
Tensor exactness & Flat-module hypothesis and exact source sequence & Tensoring multiplication by $2$ with $\Z/2\Z$ loses injectivity\\
Middle exactness & Kernel equals image & Injection, surjection, and zero composition can still leave a middle gap\\
Local differentiation & Equality on an open neighborhood & Equality only at the point does not transfer derivatives\\
Precision & Input order $N+1$ gives derivative order $N$ & Keeping the original precision after differentiation is unjustified\\
Root observation & Matching constant term and unit derivative & Wrong branch or merely nonzero nonunit derivative does not suffice\\
Behavior & Proved primitive laws and supported syntax & Passive provider-law metadata is not a Lean theorem\\
Abstraction strength & Length observer used by a length-respecting consumer & Length preservation is not sorting or extensional identity\\
Negative evidence & Concrete input violates the selected candidate's contract & Unrealizable length assignment over an empty element type is insufficient\\
Certificate shape & Exact coefficient and arity checks & Changed coefficient, wrong dimension, or a truncated multiplier list is rejected\\
Ring capability & Commutative-ring identity under the actual structure & Noncommutative input or rational witness for an integer ideal cannot be silently accepted\\
Complete solving & Soundness and coverage for the candidate predicate & A witness and coverage alone permit extraneous solutions\\
Enclosure & Collapsed bounds justify equality & A generic nondegenerate interval is not an exact value\\
Publication & All displayed claims have checked evidence & An unused false intermediate assertion still invalidates the document\\
Revision & Retained evidence rechecks at the current context & Stale results and sibling-scope proofs cannot commit edits\\
\end{longtable}}

A system that rejects both neighbors does not pass. Tests must specify the
intended rejection boundary and the expected unresolved proposition, not only
an exit code. Some weaker outputs should be accepted under a weaker request:
the correct atom-corrected CDF identity is a useful result, not a failure.

\subsection{Stopping and simplification rules}
The evaluation should permit a library-only outcome. If packages and providers
account for the improvement, retain them and drop the additional type syntax.
If property closure causes excessive search, restrict the rule set rather than
hide delays behind stronger automation. If readers confuse a refined interface
with stronger mathematics than it proves, revise the display before increasing
coverage.

Numeric success thresholds should be selected after a pilot estimates task
variation and before the main study. There is no measured productivity figure
in the present work from which to infer a credible percentage improvement.

\section{Risks, unresolved questions, and the recommended next artifact}
\label{sec:risks}
The main engineering risk is that a useful property index becomes an expensive
second elaborator with unpredictable search. The main mathematical-interface
risk is that a concise refinement hides a changed structure, scope, or
relation. The main evaluation risk is counting the benefits of new library
lemmas as benefits of syntax. The main maintenance risk is a proliferation
of redundant bundled types and transport lemmas.

Mitigations are deliberately conservative: local enrichment before packaging;
small, demand-driven rule sets; ordinary Lean contracts; selected data structures
made explicit; relation-specific adapters; and equally equipped baselines.
There is no automatic principal refinement for an arbitrary mathematical term,
and no reason to expect complete inference of all useful behavioral properties.

The next artifact should be a small Lean library and elaboration prototype,
not a broad grammar. It should prove one multivariate checker's soundness,
connect one reified expression to its original target, and expose one
ideal-preserving map composition through both ordinary Lean and refined
syntax. A paired CDF boundary example and a diagram-transport example then
test whether the same property mechanism generalizes beyond arithmetic.

A useful research question remains: how much relational structure should be
inferred from an operation's type, and how much should be supplied by named
packages? The answer should be based on failed and successful elaboration
traces. A second question is how to render a contract whose actual hypotheses
are weaker and more local than the user's familiar informal description.
The system should preserve the weaker theorem without making its statement
unreadable. Neither question requires changing the kernel first.

\section{Conclusion}
\Basalt{} treats mathematical knowledge as typed, scoped evidence about stable
objects. Unary refinements say which properties an object has; behavioral
contracts say what a selected function does; relational contracts say how an
entire diagram or collection of subjects fits together. Their composition is
ordinary mathematics, expressed by ordinary Lean theorems and used by a bounded
elaborator.

The motivating examples show why all three forms are needed. Atomlessness
licenses a CDF boundary conversion. Ideal preservation licenses quotient and
completion maps. Flatness and commuting squares license exactness transport.
A length contract licenses a particular observation of a synthesized list
program, but not arbitrary claims about its content, and an abstract
countermodel needs realizable inputs before it becomes a refutation.

This is an extension of the language in which authors state and use types,
not a proposal to weaken verification or make theorem-level equality part of
kernel conversion. The CAS and synthesis services propose evidence; they do
not choose the theorem being proved. The new surface is worth keeping only
if it improves authoring and comprehension beyond what the same theorem
packages and services already achieve in Lean.

\clearpage
\appendix
\section{Answers to the third-round questions in the two syntheses}
\label{app:questions}
The two reports overlap substantially in their questions. The following
crosswalk records concrete decisions, not a claim that every implementation
issue is settled \cite{synthesis,workbenches}.

{\small
\begin{longtable}{@{}P{29mm}P{45mm}P{71mm}@{}}
\toprule
Question family & Decision in this proposal & Required evidence or unresolved measurement\\
\midrule\endfirsthead
\toprule Question family & Decision in this proposal & Required evidence or unresolved measurement\\\midrule\endhead
\bottomrule\endfoot
First complete slice & Polynomial checker plus exact reification; ideal-preserving map wrapper & Lean soundness proof, checked execution, positive/negative cases, ordinary Lean access\\
Checker denotation & Direct evaluation into any commutative ring & Prove normalization, convolution, Boolean equality reflection, and original-target reconstruction; a bridge to semantic polynomial types is optional extra work\\
Existing tactics & Use existing proof-producing arithmetic tactics as controls & Pin and measure their actual limits; bounded failure is not nonmembership or refutation\\
Guarantee vocabulary & Ordinary propositions with typed domain-specific indices & Start with equality, guarded equality, observations, behavioral predicates, and diagram relations\\
Default semantics & Preserve imported Lean total operations & Partial-expression semantics is explicitly selected and carries its domain\\
Author delegation & Permit witness choice under a fixed requested contract & Changed domain, branch, structure, or completeness is a semantic edit\\
Visible obligations & Show restrictions that change the mathematical claim & Blind reading test determines whether default rendering is sufficient\\
Workspace state & Index Lean's actual objects and proofs, do not duplicate them & Measure new metadata cost and invalidation behavior\\
Representation routes & Start with explicitly selected theorem-backed routes & Entailment-based reuse needs checked adapters, not a universal strength score\\
First definition packages & Ideal-preserving maps and quotient observations; CDF properties as a second domain & Compare package-only and syntax-assisted authoring before adding broad schemas\\
Execution policy & Trust permission is decided before performance selection & Audit actual transitive dependencies at each version; no assumed universal speed threshold\\
Leant contribution & Structural composition plus restricted behavioral checking & Prove exact provider laws, candidate interpretation, and counterexample realizability; measure generated obligations\\
Mutation catalogue & Preserve paired tests and their distinct subcases & The present finite companion is not the implemented full-language suite\\
Publication and upgrades & Retain terms/certificates and statements at a pin & Recheck at the destination toolchain; issue a new receipt after semantic or policy changes\\
Collaboration & Merge evidence separately from changed mathematical choices & Validate contexts and selected routes after a merge; never simply union stale caches\\
Stopping rule & Library, adapter, or renderer alone is an acceptable result & New syntax must add measurable benefit with shared infrastructure and acceptable semantic-reading accuracy\\
\end{longtable}}

The reports' invitation to execute a real Lean slice is appropriate. This
article does not claim to have met that gate: no Lean compiler was available
in the working environment. It instead supplies a more explicit mathematical
checker specification, a small uncompiled proof-term core, and executable
finite demonstrations with their scope stated. A kernel-checked prototype is
the next concrete deliverable, not an accomplishment inferred from this report.

\section{A minimal ordinary Lean core, with its validation status}
\label{app:lean}
The companion \code{ContractCore.lean} illustrates the logical encoding without
new syntax, external solvers, or a dependency on Mathlib. It contains no
\code{sorry} or added axiom. It was \emph{not compiled here}; absence of those
strings is not evidence that Lean accepted it. The following excerpt is
original illustrative code, not copied from a repository.

\par\noindent\begin{minipage}{\linewidth}
\begin{lstlisting}
structure LawfulFn (A : Type u) (B : Type v)
    (pre : A -> Prop) (post : A -> B -> Prop) where
  run : A -> B
  law : forall x, pre x -> post x (run x)

def composeLaw
    (f : A -> B) (g : B -> C)
    (P : A -> Prop) (Q : A -> B -> Prop)
    (R : A -> B -> C -> Prop)
    (hf : forall x, P x -> Q x (f x))
    (hg : forall x y, Q x y -> R x y (g y)) :
    forall x, P x -> R x (f x) (g (f x)) :=
  fun x hx => hg x (f x) (hf x hx)
\end{lstlisting}
\end{minipage}\par

This core also includes value-preserving refinement weakening, intersection
packaging, the contract consequence rule, and observation-based congruence.
These definitions illustrate why the type-system proposal can remain
conservative. They are not an implementation of the fact index, parser,
provider protocol, or polynomial checker in Lean.

\section{Source register and scope of inspection}
\label{app:sources}
All sources were accessed on 6 September 2026. Repository heads were resolved
through the GitHub connector; the following commits identify the observed
snapshots. Individual blob identities provide finer provenance for the files
actually used. No repository was modified.

{\small
\begin{tabularx}{\textwidth}{@{}P{34mm}Y@{}}
\toprule
Repository & Observed commit\\
\midrule
Brainstorm & \code{b052ec011beb96a7c48df2a9988823f0f779b445}\\
ProveIt & \code{24ce8bd743eaab64a91ce90725ea00f498d319d2}\\
Leant & \code{3a40904be8a410d832d9ab6900b3d3e7b425eccb}\\
FLT & \code{aa2d8b34692b16c70f699536de0d8e75b9a3e9ef}\\
\bottomrule
\end{tabularx}}

\subsection{Both requested syntheses}
\code{docs/round-2/synthesis/unified-report.tex}, blob
\code{ae7c21dac1ba615e36dff07ea6f6a1c433a40108}, was read through its
substantive sections, corrections, questions, appendices, and bibliography.
\code{docs/round-2/unified_report/unified_report.tex}, blob
\code{de10daae20af19285ee07a476caa94e577f7d7d9}, was likewise read closely.
The original nine proposals are discussed through these syntheses; this work
does not claim an independent full reading of all nine reports or reproduction
of their source inventories and experiments.

\subsection{ProveIt source modules}
\code{Analysis/FabiusFunction/Lean/FabiusFunction/ContinuousCDF.lean}, blob
\code{2306bd53bd46d145b4e156c065b058009ce16fdc}, was read in full. The
relevant declarations are \code{continuous_cdf_of_nullSingleton},
\code{cdf_reflection_sub}, and
\code{measure_eq_withDensity_of_cdf_hasDerivAt}.

\code{Analysis/FabiusFunction/Lean/FabiusFunction/AutonomousIteratedDeriv.lean},
blob \code{1422b87f31073fed898ff0c8a080ea0a798aef5e}, was read in full.
The main interface is \code{AutonomousODE.iteratedDeriv_eq_comp}, with its
more globally differentiable corollary used only as a comparison.

\subsection{FLT source modules}
\code{Definitions/Def_AdicCompletionRingFunctoriality.lean}, blob
\code{13acc67dd2b0114d0f1189584b847748e6e59654}, was read in full. The
relevant interfaces include \code{levelMap} and completion-map constructions
(the source uses subscripted algebra-map names), their evaluation and
composition lemmas, \code{map_comp_le}, and the equivalence constructions
from compatible levelwise bijections.

\code{P2M/Sol/S_Rep_indBotSC_shortExact.lean}, blob
\code{63caa90be75bb6875f1211a8e81230f00931fe5f}, was read in full.
The relevant parts are the generator-based modeling equivalence, its
\code{e_square} naturality theorem, \code{shortExact_of_ladder}, and the
final \code{solution} theorem. The repository README, blob
\code{f3cfcb92443c29a8ce87494f6c9a7b57e40129e9}, was also read.
Its complete-build and independent-checking claims are attributed to that
source, not reproduced results of this work.

\subsection{Leant source modules}
\code{src/Leant/Synth/Verification.hs}, blob
\code{caee3c51cfd226e6c2bb9887f80293343abd486c}, and
\code{src/Leant/Synth/Length/Contract.hs}, blob
\code{a1f04001c6ccd6c5e018162fa4db7ff5b059972e}, were read in full.

The first 200 source lines of \code{src/Leant/Synth/Engine.hs}, blob
\code{1064c8215506e172401cd7c5c06d7e8222e0befa}, and the first 220 lines of
\code{src/Leant/Synth/Length/Handoff.hs}, blob
\code{883e6fd2147fe9aa100c1691fb7fdb12b0f9380d}, were inspected.
They support the discussion of exported typed origins, explicit handoff
contracts, retained provenance, and refusal conditions. This is not a
whole-engine correctness audit or a solver execution test.

\subsection{Companion validation and reproducibility}
The archive contains the self-contained TeX source and built PDF, a build
script, the finite Python companion, its actual test output and JSON receipt,
and the uncompiled Lean core. The Python companion uses only the standard
library. Its 24 passing test methods include several loops over small finite
inputs; the method count is not a count of universal theorems or implemented
language features.

The PDF build is checked for unresolved references and layout warnings and
rendered for visual inspection. These document checks say nothing about
mathematical kernel acceptance. The output manifest records hashes of the
final distributable files; it does not claim to hash independently downloaded
copies of the upstream repositories.

\clearpage
\begin{thebibliography}{99}
\small\raggedright
\setlength{\parskip}{0pt}
\setlength{\itemsep}{2pt}
\bibitem{synthesis} Codex. \emph{Mathematical Objects, Certified Computation:
A Unified Evaluation of Nine Proof-Language Proposals}. Round-two synthesis,
revised 6 September 2026. Brainstorm repository, source and blob listed in
Appendix~\ref{app:sources}.
\url{https://github.com/VladimirReshetnikov/Brainstorm/blob/b052ec011beb96a7c48df2a9988823f0f779b445/docs/round-2/synthesis/unified-report.tex}.

\bibitem{workbenches} Claude. \emph{Nine Workbenches: Certified Computation in a
Mathematician's Proof Language over Lean}. Round-two unified review, title
dated 5 September 2026, revised version inspected 6 September 2026.
\url{https://github.com/VladimirReshetnikov/Brainstorm/blob/b052ec011beb96a7c48df2a9988823f0f779b445/docs/round-2/unified_report/unified_report.tex}.

\bibitem{cdfsource} V. Reshetnikov and contributors. ProveIt,
\emph{ContinuousCDF.lean}. Source inspected at the registered revision.
\url{https://github.com/VladimirReshetnikov/ProveIt/blob/24ce8bd743eaab64a91ce90725ea00f498d319d2/Analysis/FabiusFunction/Lean/FabiusFunction/ContinuousCDF.lean}.

\bibitem{odesource} V. Reshetnikov and contributors. ProveIt,
\emph{AutonomousIteratedDeriv.lean}.
\url{https://github.com/VladimirReshetnikov/ProveIt/blob/24ce8bd743eaab64a91ce90725ea00f498d319d2/Analysis/FabiusFunction/Lean/FabiusFunction/AutonomousIteratedDeriv.lean}.

\bibitem{fltreadme} Anthropic and credited contributors. \emph{Fermat's Last
Theorem in Lean 4}, README and source attribution.
\url{https://github.com/anthropics/fermats-last-theorem/blob/aa2d8b34692b16c70f699536de0d8e75b9a3e9ef/README.md}.

\bibitem{completionsource} Anthropic and credited contributors.
\emph{Def\_AdicCompletionRingFunctoriality.lean}.
\url{https://github.com/anthropics/fermats-last-theorem/blob/aa2d8b34692b16c70f699536de0d8e75b9a3e9ef/Definitions/Def_AdicCompletionRingFunctoriality.lean}.

\bibitem{exactsource} Anthropic and credited contributors.
\emph{S\_Rep\_indBotSC\_shortExact.lean}.
\url{https://github.com/anthropics/fermats-last-theorem/blob/aa2d8b34692b16c70f699536de0d8e75b9a3e9ef/P2M/Sol/S_Rep_indBotSC_shortExact.lean}.

\bibitem{leantverification} V. Reshetnikov and contributors. Leant,
\emph{Leant.Synth.Verification}.
\url{https://github.com/VladimirReshetnikov/Leant/blob/3a40904be8a410d832d9ab6900b3d3e7b425eccb/src/Leant/Synth/Verification.hs}.

\bibitem{leantengine} V. Reshetnikov and contributors. Leant,
\emph{Leant.Synth.Engine}, header and exported interface inspected.
\url{https://github.com/VladimirReshetnikov/Leant/blob/3a40904be8a410d832d9ab6900b3d3e7b425eccb/src/Leant/Synth/Engine.hs}.

\bibitem{leantcontract} V. Reshetnikov and contributors. Leant,
\emph{Leant.Synth.Length.Contract}.
\url{https://github.com/VladimirReshetnikov/Leant/blob/3a40904be8a410d832d9ab6900b3d3e7b425eccb/src/Leant/Synth/Length/Contract.hs}.

\bibitem{leanthandoff} V. Reshetnikov and contributors. Leant,
\emph{Leant.Synth.Length.Handoff}, first 220 lines inspected.
\url{https://github.com/VladimirReshetnikov/Leant/blob/3a40904be8a410d832d9ab6900b3d3e7b425eccb/src/Leant/Synth/Length/Handoff.hs}.

\bibitem{leansubtypes} Lean team. \emph{Subtypes}. Lean Language Reference,
latest documentation accessed 6 September 2026.
\url{https://lean-lang.org/doc/reference/latest/Basic-Types/Subtypes/}.

\bibitem{leanaxioms} Lean team. \emph{Axioms}. Lean Language Reference,
including the discussion of nonparametric polymorphic functions. Accessed
6 September 2026.
\url{https://lean-lang.org/doc/reference/latest/Axioms/}.

\bibitem{leanvalidation} Lean team. \emph{Validating a Lean Proof}. Lean
Language Reference. Accessed 6 September 2026.
\url{https://lean-lang.org/doc/reference/latest/ValidatingProofs/}.

\bibitem{lovaspfenning} William Lovas and Frank Pfenning.
\emph{Refinement Types for Logical Frameworks and Their Interpretation as
Proof Irrelevance}. 2010. arXiv:1009.1861. Abstract and publication metadata
consulted; used here as a precedent, not as a technical dependency.
\url{https://arxiv.org/abs/1009.1861}.

\bibitem{explicitrefinements} Jad Elkhaleq Ghalayini and Neel Krishnaswami.
\emph{Explicit Refinement Types}. 2023. arXiv:2311.13995. Abstract and
publication metadata consulted.
\url{https://arxiv.org/abs/2311.13995}.
\end{thebibliography}
\end{document}
